% Auto-generated from section_b_final_conclusions_report.md.
% Re-run tools/convert_section_b_md_to_tex.py after editing the Markdown source.
\documentclass[UTF8,12pt]{ctexart}

\usepackage[a4paper,margin=2.2cm]{geometry}
\usepackage{amsmath,amssymb}
\usepackage{array}
\usepackage{booktabs}
\usepackage{caption}
\usepackage{enumitem}
\usepackage{float}
\usepackage{fvextra}
\usepackage{graphicx}
\usepackage{longtable}
\usepackage{newunicodechar}
\usepackage{seqsplit}
\usepackage{titlesec}
\usepackage{xcolor}
\usepackage{hyperref}

\graphicspath{{./}}
\hypersetup{
  colorlinks=true,
  linkcolor=blue!55!black,
  urlcolor=blue!55!black,
  citecolor=blue!55!black
}
\setlength{\parindent}{2em}
\setlength{\parskip}{0.35em}
\setlength{\emergencystretch}{3em}
\sloppy
\setcounter{tocdepth}{3}
\setlist{leftmargin=2em,itemsep=0.15em,topsep=0.25em}
\renewcommand{\arraystretch}{1.22}
\titleformat{\paragraph}[block]{\normalfont\normalsize\bfseries}{}{0pt}{}
\titlespacing*{\paragraph}{0pt}{1.1ex plus .2ex}{0.6ex}
\DefineVerbatimEnvironment{CodeBlock}{Verbatim}{breaklines=true,breakanywhere=true,fontsize=\small}
\newcommand{\CodeInline}[1]{\texttt{\seqsplit{#1}}}
\newunicodechar{≥}{\ensuremath{\ge}}
\newunicodechar{≤}{\ensuremath{\le}}
\newunicodechar{≈}{\ensuremath{\approx}}
\newunicodechar{∈}{\ensuremath{\in}}
\newunicodechar{∧}{\ensuremath{\land}}
\newunicodechar{ρ}{\ensuremath{\rho}}
\newunicodechar{τ}{\ensuremath{\tau}}
\newunicodechar{σ}{\ensuremath{\sigma}}
\newunicodechar{✓}{\ensuremath{\checkmark}}
\newunicodechar{→}{\ensuremath{\to}}

\title{CET-6 Section B 答案位置规律分析}
\author{}
\date{\today}

\begin{document}
\maketitle
\tableofcontents
\newpage

\phantomsection
\section*{摘要}
\addcontentsline{toc}{section}{摘要}


这份报告使用已经人工核对的 CET-6 Section B 数据，考察 \CodeInline{Q36-Q45} 的段落位置、答案字母、题号之间的无条件关系、条件联动、整套容量约束以及这些统计量在考场中的用法。正文先说明样本口径、变量和图表读法，再讨论条件窗口、边缘带、镜像关系和各题的实际角色；统计结果与做题建议分开陈述，避免把相关性误当成规则。


正文的统计结论集中在第 1-14 章，第 15-21 章保留核对记录、证据表和操作手册。附录中出现的“信息增益最大”“残余熵最低”等说法只对应相应指标；考试中的先后顺序仍以第 14.1 节的口径为准。


这份报告得到的结果并不能由单题均值概括。均值只给出一个大概方向，实际判断还要把单题分布、整套容量、边缘带、条件联动和反向排除放在一起看。题号相邻也不意味着原文位置相邻；在这批样本中，隔一题的同奇偶链和少数高杠杆枢纽题反而更值得联动。\CodeInline{Q36} 更像一个位置锚，而不是默认的起手锚；\CodeInline{Q44} 虽然不是尾题，却承担着中后段的联动作用；\CodeInline{Q38} 则属于“做起来不稳定，做出后很有反推价值”的题。

位置落在 \CodeInline{0-20\%} 或 \CodeInline{80-100\%} 这类边缘带时，位置本身就带有比均值更多的条件信息。因而考场上的用法不是猜某题会落在哪个平均位置，而是先用粗分布确定大区，再用容量约束检查某一段是不是已经挤满，最后才用条件规则收窄搜索窗。下面的统计和策略，都是围绕这个先后关系展开的；if-else 只是把已经形成的判断顺序固定下来。


\phantomsection
\section*{1. 引言}
\addcontentsline{toc}{section}{1. 引言}


\phantomsection
\subsection*{1.1 研究问题}
\addcontentsline{toc}{subsection}{1.1 研究问题}


面对 Section B 这类信息匹配题，常见的直觉有两种：有些题天生偏前或偏后，相邻题也应当在原文中相邻。第一种看法顶多给出方向，第二种看法在这批数据里并不可靠。真正需要回答的不是“\CodeInline{Q41} 的平均位置是不是更靠后”，而是：哪些题适合先做，哪一道题做出后最能带动其他题，某题落到极前或极后时会把哪些候选推开，以及一个 20\% 区间已经装下 3 题后是否还值得继续往里找。后文的统计，正是为了把这些看似零散的考场判断接到同一套证据上。


\phantomsection
\subsection*{1.2 报告结构}
\addcontentsline{toc}{subsection}{1.2 报告结构}


阅读时可以先沿主报告往下看：先交代数据和样本口径，再说明变量、文件和图表的读法，随后进入单题位置、答案字母、题号角色、题间关系和条件联动，最后才把统计结果翻译成考场策略，并说明它的局限。第 1--14 章负责这条主线，第 15--21 章保留复核过程、速记版、考试手册和深度分析结果；第 18--20 章是证据附录，不另起一套策略，也不覆盖第 12--14 章和第 17 章的结论。


\phantomsection
\subsection*{1.3 文档合并说明}
\addcontentsline{toc}{subsection}{1.3 文档合并说明}


为便于查找，本文作为 \CodeInline{\_analysis\_output} 下 Section B 分析的\textbf{单一主入口}。\CodeInline{section\_b\_correlation\_notes.md} 负责题间关系，\CodeInline{sample\_answer\_second\_pass\_report.md} 负责二次核对，\CodeInline{section\_b\_exam\_playbook.md} 负责考场执行；这三份结论型 Markdown 的内容已经并入正文。


只想看最终结果时，读这一个文件即可。


\CodeInline{verification\_packets/}、\CodeInline{recheck\_packets/} 和 \CodeInline{recheck\_reviews/} 下的 Markdown 是逐篇核对包、批次复核记录和工作底稿。它们不再承担“主结论报告”的角色，附录只说明各自用途和入口。


\phantomsection
\section*{2. 数据与样本口径}
\addcontentsline{toc}{section}{2. 数据与样本口径}


\phantomsection
\subsection*{2.1 当前样本口径已经统一}
\addcontentsline{toc}{subsection}{2.1 当前样本口径已经统一}


这一版中，答案层和位置层使用同一批样本，不再需要把“答案字母统计”和“段落位置统计”分成两套口径。


\CodeInline{manual\_truth\_summary.json} 中的四个计数分别是 \CodeInline{files\_count = 59}、\CodeInline{titles\_count = 59}、\CodeInline{records\_count = 590} 和 \CodeInline{source\_type\_counts.manual = 590}。换句话说，当前报告使用的是 59 套文章、每套 10 题的完整人工核对结果。


位置统计则来自 \CodeInline{question\_position\_stats.csv}、\CodeInline{question\_pair\_relation\_stats.csv}、\CodeInline{question\_order\_profile.csv} 和 \CodeInline{conditional\_window\_gain\_stats.csv}，这些文件都以同一批文章中的 \CodeInline{paragraph\_position\_pct} 为基础。


当前总整理口径和位置分析口径都是 \CodeInline{59} 套 /\allowbreak{} \CodeInline{590} 题。因此，段落位置、相关性、容量约束和条件联动使用的是同一分母，不需要在不同章节之间来回换算。


\phantomsection
\subsection*{2.2 历史上为何曾经出现两层口径}
\addcontentsline{toc}{subsection}{2.2 历史上为何曾经出现两层口径}


早期分析曾出现两层口径，原因是有 \CodeInline{3} 篇 2024 年 6 月文章只有人工复核的答案字母，没有可靠的标准化段落位置百分比。


当时答案层已经做到 \CodeInline{59} 套 /\allowbreak{} \CodeInline{590} 题，位置层却只能做到 \CodeInline{56} 套 /\allowbreak{} \CodeInline{560} 题，所以答案字母分布与段落位置、题间距离、条件联动不得不共存两套分母。


后来补入 \CodeInline{section\_b\_texts} 标准链路的 3 篇文章是 \CodeInline{The Curious Case of the Tree That Owns Itself}、\CodeInline{These are the habits to avoid if you want to make a behavior change} 和 \CodeInline{Blame your worthless workdays on meeting recovery syndrome}。


补齐后，当前版本不再有“双口径”。后文若出现旧草稿中的 \CodeInline{56}、\CodeInline{560} 或 \CodeInline{ocr\_manual}，它们只是 \textbf{旧阶段工作痕迹}，不代表现在的样本范围。


\phantomsection
\subsection*{2.3 当前数据适合回答什么}
\addcontentsline{toc}{subsection}{2.3 当前数据适合回答什么}


当前样本并不算大，但已经足以判断前、中、后段的大致分布，比较相邻题与隔一题链的强弱，并挖出一批可在考场使用的条件联动规则。它还不足以给出逐题逐段的高置信硬预测，也不足以让每一条小样本条件规则都承担强结论；这些边界会在后文的解释中保留。


\phantomsection
\section*{3. 变量、文件与图表说明}
\addcontentsline{toc}{section}{3. 变量、文件与图表说明}


\phantomsection
\subsection*{3.1 核心字段说明}
\addcontentsline{toc}{subsection}{3.1 核心字段说明}


后文反复使用的字段如下：


\begin{center}
\small
\setlength{\tabcolsep}{4pt}
\begin{longtable}{@{}>{\raggedright\arraybackslash}p{0.460\linewidth}>{\raggedright\arraybackslash}p{0.460\linewidth}@{}}
\toprule
\textbf{字段} & \textbf{含义} \\
\midrule
\endfirsthead
\toprule
\textbf{字段} & \textbf{含义} \\
\midrule
\endhead
\CodeInline{question\_number} & 题号，范围 \CodeInline{36-45} \\
\CodeInline{answer\_letter} & 该题对应原文段落字母，不是选择题选项 \\
\CodeInline{paragraph\_position\_pct} & 该题对应段落在全文中的相对位置百分比 \\
\CodeInline{paragraph\_count} & 原文总段落数 \\
\CodeInline{dominant\_band} & 最常落入的 20\% 位置带 \\
\CodeInline{dominant\_10\_band} & 最常落入的 10\% 位置带 \\
\CodeInline{band\_entropy\_20} & 20\% 分带下的分散度，越高越分散 \\
\CodeInline{band\_entropy\_10} & 10\% 分带下的分散度，越高越分散 \\
\CodeInline{spearman\_pct} & 两题相对位置的秩相关 \\
\CodeInline{within\_2\_paragraph\_pct} & 两题段落差绝对值不超过 2 的比例 \\
\CodeInline{same\_half\_pct} & 两题同在前半区或后半区的比例 \\
\CodeInline{median\_gap\_paragraphs} & 两题段落差的中位数 \\
\CodeInline{best\_contiguous\_40\_coverage\_pct} & 某题最佳 40\% 连续窗口覆盖率 \\
\CodeInline{gain\_40\_pct} & 条件成立后，对目标题 40\% 搜索窗覆盖率的提升 \\
\CodeInline{rule\_count\_gain15} & 某题作为条件源时，能带来至少 15\% 窗口增益的规则数 \\
\bottomrule
\end{longtable}
\normalsize
\end{center}


\phantomsection
\subsection*{3.2 关键输出文件说明}
\addcontentsline{toc}{subsection}{3.2 关键输出文件说明}


\begin{center}
\small
\setlength{\tabcolsep}{4pt}
\begin{longtable}{@{}>{\raggedright\arraybackslash}p{0.460\linewidth}>{\raggedright\arraybackslash}p{0.460\linewidth}@{}}
\toprule
\textbf{文件} & \textbf{作用} \\
\midrule
\endfirsthead
\toprule
\textbf{文件} & \textbf{作用} \\
\midrule
\endhead
\CodeInline{manual\_truth\_records.csv} & 基础真值表 \\
\CodeInline{question\_position\_stats.csv} & 单题位置总体统计 \\
\CodeInline{question\_position\_fine\_stats.csv} & 单题 10\% 热区统计 \\
\CodeInline{question\_answer\_summary.csv} & 单题答案字母分布摘要 \\
\CodeInline{letter\_to\_question\_summary.csv} & 反向字母到题号分布摘要 \\
\CodeInline{question\_pair\_relation\_stats.csv} & 任意两题之间的无条件关系 \\
\CodeInline{adjacent\_question\_relation\_stats.csv} & 只看相邻题的关系摘要 \\
\CodeInline{question\_strategy\_features.csv} & 单题策略特征汇总 \\
\CodeInline{question\_order\_profile.csv} & 每题在整套题中的先后 rank 分布 \\
\CodeInline{conditional\_window\_gain\_stats.csv} & 条件联动规则总表 \\
\CodeInline{anchor\_leverage\_summary.csv} & 高杠杆题汇总 \\
\CodeInline{question\_capacity\_summary.csv} & 容量约束摘要 \\
\bottomrule
\end{longtable}
\normalsize
\end{center}


\phantomsection
\subsection*{3.3 图表怎么读，怎么用}
\addcontentsline{toc}{subsection}{3.3 图表怎么读，怎么用}


\phantomsection
\subsubsection*{\CodeInline{question\_position\_boxplot.png}}
\addcontentsline{toc}{subsubsection}{question\_position\_boxplot.png}
\begin{figure}[H]
\centering
\includegraphics[width=0.92\linewidth]{question_position_boxplot.png}
\caption{question\_position\_boxplot.png}
\end{figure}


箱线图先回答位置大致偏前还是偏后，再回答这种倾向稳不稳定。中位数给出中心位置，四分位距反映主要样本散得多开，两端的长尾则提醒我们还有没有少数反常落点。只有中位数靠前且箱体较窄时，“偏前”才是较稳定的描述；若箱体很宽或长尾明显，同一个中位数就不能被拿来押具体段落。


\phantomsection
\subsubsection*{\CodeInline{question\_position\_band\_distribution.png}}
\addcontentsline{toc}{subsubsection}{question\_position\_band\_distribution.png}
\begin{figure}[H]
\centering
\includegraphics[width=0.92\linewidth]{question_position_band_distribution.png}
\caption{question\_position\_band\_distribution.png}
\end{figure}


分带图把全文切成五个 20\% 大区间，适合先判断某道题的候选主要堆在哪一片。考试时可以据此决定从前、中、后哪个方向开始找，却不能继续把一个大带解释成某个具体段落；它解决的是搜索方向，不是最终落点。


\phantomsection
\subsubsection*{\CodeInline{question\_position\_band\_heatmap.png}}
\addcontentsline{toc}{subsubsection}{question\_position\_band\_heatmap.png}
\begin{figure}[H]
\centering
\includegraphics[width=0.92\linewidth]{question_position_band_heatmap.png}
\caption{question\_position\_band\_heatmap.png}
\end{figure}


热力图继续把位置细分到 10\% 区间。这里更值得看的不是最深的那一格本身，而是深色区域连成一片、分成两峰，还是散落在多个区间：前一种说明搜索窗较集中，双峰意味着需要保留两条路线，散峰则说明题号本身提供不了稳定定位。


\phantomsection
\subsubsection*{\CodeInline{question\_answer\_distribution.png}}
\addcontentsline{toc}{subsubsection}{question\_answer\_distribution.png}
\begin{figure}[H]
\centering
\includegraphics[width=0.92\linewidth]{question_answer_distribution.png}
\caption{question\_answer\_distribution.png}
\end{figure}


答案分布图直接显示每道题在原文段落字母上的历史落点。高频字母可以提示前后段倾向，但字母还受到文章段落数差异的影响，所以它更适合作为已有判断的补充：题干已经把候选压到几段时，高频位置可以帮助排序；候选尚未形成时，单凭一个高频字母下注并不稳妥。


\phantomsection
\subsubsection*{\CodeInline{letter\_to\_question\_heatmap.png}}
\addcontentsline{toc}{subsubsection}{letter\_to\_question\_heatmap.png}
\begin{figure}[H]
\centering
\includegraphics[width=0.92\linewidth]{letter_to_question_heatmap.png}
\caption{letter\_to\_question\_heatmap.png}
\end{figure}


反向热力图换了一个问题：当某个段落字母已经进入候选时，历史上更常由哪些题号占用它。这个方向在后期排除中更有用，例如同一字母同时被几题争夺时，可以用它检查哪种分配更少见；若还没有读题干、候选也没有缩小，反向频率本身不能替代内容匹配。


几张图放在一起后，各自的作用并不相同：箱线图判断稳定程度，分带图确定搜索方向，细粒度热力图保留局部峰值，字母图在候选收窄后参与排除。它们共同完成粗筛和联动，但没有任何一张图能够脱离题干单独定点。


\phantomsection
\subsection*{3.4 图表文件索引}
\addcontentsline{toc}{subsection}{3.4 图表文件索引}


正文使用的图表文件均在 \CodeInline{cet 6/\_analysis\_output/}：


\begin{itemize}
\item \CodeInline{question\_position\_boxplot.png}：单题位置箱线图，用来看中位数、四分位距与离散度。
\item \CodeInline{question\_position\_band\_distribution.png}：单题在五个 \CodeInline{20\%} 大带上的分布图，用来做粗定位。
\item \CodeInline{question\_position\_band\_heatmap.png}：单题在 \CodeInline{10\%} 细带上的热区图，用来观察双峰、散峰和局部集中。
\item \CodeInline{question\_position\_curve.png}：单题位置曲线图，用来辅助观察整体形态、长尾与多峰倾向。
\item \CodeInline{question\_answer\_distribution.png}：题号到原文段落字母的分布图，用来做后期排除与偏向判断。
\item \CodeInline{letter\_to\_question\_heatmap.png}：字母反向映射到题号的热力图，用来辅助判断某字母更常落在哪些题号上。
\end{itemize}


\phantomsection
\section*{4. 单题位置分布与离散程度}
\addcontentsline{toc}{section}{4. 单题位置分布与离散程度}


\phantomsection
\subsection*{4.1 粗粒度 20\% 分带分布}
\addcontentsline{toc}{subsection}{4.1 粗粒度 20\% 分带分布}


根据 \CodeInline{question\_position\_stats.csv}：


\begin{center}
\small
\setlength{\tabcolsep}{4pt}
\begin{longtable}{@{}>{\raggedright\arraybackslash}p{0.172\linewidth}>{\raggedright\arraybackslash}p{0.172\linewidth}>{\raggedright\arraybackslash}p{0.172\linewidth}>{\raggedright\arraybackslash}p{0.172\linewidth}>{\raggedright\arraybackslash}p{0.172\linewidth}@{}}
\toprule
\textbf{题号} & \textbf{主导 20\% 带} & \textbf{占比} & \textbf{次主导 20\% 带} & \textbf{占比} \\
\midrule
\endfirsthead
\toprule
\textbf{题号} & \textbf{主导 20\% 带} & \textbf{占比} & \textbf{次主导 20\% 带} & \textbf{占比} \\
\midrule
\endhead
Q36 & \CodeInline{20-40\%} & 37.50\% & \CodeInline{40-60\%} & 33.93\% \\
Q37 & \CodeInline{0-20\%} & 30.36\% & \CodeInline{60-80\%} & 26.79\% \\
Q38 & \CodeInline{0-20\%} & 26.79\% & \CodeInline{60-80\%} & 23.21\% \\
Q39 & \CodeInline{20-40\%} & 39.29\% & \CodeInline{40-60\%} & 21.43\% \\
Q40 & \CodeInline{40-60\%} & 28.57\% & \CodeInline{20-40\%} & 21.43\% \\
Q41 & \CodeInline{80-100\%} & 32.14\% & \CodeInline{40-60\%} & 28.57\% \\
Q42 & \CodeInline{0-20\%} & 26.79\% & \CodeInline{20-40\%} & 21.43\% \\
Q43 & \CodeInline{40-60\%} & 28.57\% & \CodeInline{0-20\%} & 21.43\% \\
Q44 & \CodeInline{20-40\%} & 30.36\% & \CodeInline{0-20\%} & 21.43\% \\
Q45 & \CodeInline{60-80\%} & 32.14\% & \CodeInline{80-100\%} & 23.21\% \\
\bottomrule
\end{longtable}
\normalsize
\end{center}


主导分带把十道题分成了三种不同角色。\CodeInline{Q36 / Q39 / Q44} 的高频区域主要在前中段，\CodeInline{Q41 / Q45} 明显向后移动；\CodeInline{Q38 / Q42 / Q43} 的主导带占比都不高，单凭大区分布很难把搜索范围压窄。


\phantomsection
\subsection*{4.2 细粒度 10\% 热区分布}
\addcontentsline{toc}{subsection}{4.2 细粒度 10\% 热区分布}


根据 \CodeInline{question\_position\_fine\_stats.csv}：


\begin{center}
\small
\setlength{\tabcolsep}{4pt}
\begin{longtable}{@{}>{\raggedright\arraybackslash}p{0.172\linewidth}>{\raggedright\arraybackslash}p{0.172\linewidth}>{\raggedright\arraybackslash}p{0.172\linewidth}>{\raggedright\arraybackslash}p{0.172\linewidth}>{\raggedright\arraybackslash}p{0.172\linewidth}@{}}
\toprule
\textbf{题号} & \textbf{主导 10\% 带} & \textbf{占比} & \textbf{次主导 10\% 带} & \textbf{占比} \\
\midrule
\endfirsthead
\toprule
\textbf{题号} & \textbf{主导 10\% 带} & \textbf{占比} & \textbf{次主导 10\% 带} & \textbf{占比} \\
\midrule
\endhead
Q36 & \CodeInline{40-50\%} & 25.00\% & \CodeInline{20-30\%} & 19.64\% \\
Q37 & \CodeInline{10-20\%} & 25.00\% & \CodeInline{70-80\%} & 17.86\% \\
Q38 & \CodeInline{20-30\%} & 16.07\% & \CodeInline{0-10\%} & 14.29\% \\
Q39 & \CodeInline{20-30\%} & 21.43\% & \CodeInline{30-40\%} & 17.86\% \\
Q40 & \CodeInline{40-50\%} & 19.64\% & \CodeInline{20-30\%} & 12.50\% \\
Q41 & \CodeInline{40-50\%} & 21.43\% & \CodeInline{80-90\%} & 19.64\% \\
Q42 & \CodeInline{10-20\%} & 14.29\% & \CodeInline{0-10\%} & 12.50\% \\
Q43 & \CodeInline{50-60\%} & 16.07\% & \CodeInline{70-80\%} & 14.29\% \\
Q44 & \CodeInline{30-40\%} & 17.86\% & \CodeInline{20-30\%} & 12.50\% \\
Q45 & \CodeInline{60-70\%} & 21.43\% & \CodeInline{80-90\%} & 17.86\% \\
\bottomrule
\end{longtable}
\normalsize
\end{center}


细分到 10\% 后，\CodeInline{Q37} 在 \CodeInline{10-20\%} 和 \CodeInline{70-80\%} 各有一处高发区，前后双峰比大带表呈现得更直接。\CodeInline{Q41} 的峰值分布在 \CodeInline{40-50\%} 与 \CodeInline{80-90\%}，也不能只保留一个后段落点；\CodeInline{Q45} 的两个高频区都在后半段，方向相对一致。相比之下，\CodeInline{Q38 / Q42 / Q43 / Q44} 的最高占比仍不突出，细分并没有消除其位置摆动。


\phantomsection
\subsection*{4.3 分散度指标说明}
\addcontentsline{toc}{subsection}{4.3 分散度指标说明}


\CodeInline{band\_entropy\_20} 和 \CodeInline{band\_entropy\_10} 越高，说明该题越分散。


三道高熵题的数值很接近：\CodeInline{Q38} 的 \CodeInline{band\_entropy\_20 = 0.9849}、\CodeInline{band\_entropy\_10 = 0.9603}；\CodeInline{Q42} 分别为 \CodeInline{0.9768} 和 \CodeInline{0.9792}；\CodeInline{Q44} 分别为 \CodeInline{0.9663} 和 \CodeInline{0.9713}。粗分带和细分带下都保持较高熵值，说明分散不是某一种划分方式造成的。它们更适合在其他题提供条件后再处理。

作为对照，\CodeInline{Q36} 的 \CodeInline{band\_entropy\_20 = 0.8606}，位置集中程度高于上述几题。它能够帮助判断大区，但这并不等于它必须最先作答；\CodeInline{Q38 / Q42 / Q43} 则更依赖已形成的条件窗口。


\phantomsection
\subsection*{4.4 单题连续窗口覆盖率}
\addcontentsline{toc}{subsection}{4.4 单题连续窗口覆盖率}


\CodeInline{question\_strategy\_features.csv} 给出了不附加其他题条件时的最佳连续窗口。\CodeInline{Q36} 与 \CodeInline{Q39} 的窗口都是 \CodeInline{20-60\%}，覆盖率分别为 \CodeInline{71.43\%} 和 \CodeInline{60.71\%}；\CodeInline{Q41} 与 \CodeInline{Q45} 的窗口都移到 \CodeInline{60-100\%}，覆盖率分别为 \CodeInline{44.64\%} 和 \CodeInline{55.36\%}。这些窗口适合确定裸做时的阅读大区，覆盖率本身也提醒我们，后两题仍需依靠题干和联动继续缩小范围。


\phantomsection
\section*{5. 答案字母分布层面还能看到什么}
\addcontentsline{toc}{section}{5. 答案字母分布层面还能看到什么}


\phantomsection
\subsection*{5.1 单题对应字母分布}
\addcontentsline{toc}{subsection}{5.1 单题对应字母分布}


\CodeInline{question\_answer\_summary.csv} 中，\CodeInline{Q36} 最常见的字母是 \CodeInline{E}，占 \CodeInline{20.34\%}；\CodeInline{Q37} 最常见的是 \CodeInline{C}，占 \CodeInline{22.03\%}。即便是各自最高的字母，占比也只在两成左右。\CodeInline{Q41} 和 \CodeInline{Q44} 的前两高字母还出现并列，\CodeInline{top\_gap\_pct = 0}。主峰不尖意味着字母只能给已有候选增加一点权重，无法单独支撑押题。


\phantomsection
\subsection*{5.2 反向看：某字母更偏向哪些题号}
\addcontentsline{toc}{subsection}{5.2 反向看：某字母更偏向哪些题号}


反向统计得到的集中程度略有提高。\CodeInline{N} 最常见于 \CodeInline{Q45}，占 \CodeInline{40.0\%}；\CodeInline{C} 最常见于 \CodeInline{Q37}，占 \CodeInline{27.66\%}；\CodeInline{E} 和 \CodeInline{F} 最常见于 \CodeInline{Q36}，占比分别为 \CodeInline{24.49\%} 和 \CodeInline{24.44\%}。不过，大多数字母的 \CodeInline{reverse\_entropy} 仍不低。只有当某个字母已经进入少量候选时，这些比例才适合参与排除；它们不能取代题干内容成为主判断依据。


\phantomsection
\subsection*{5.3 字母层面的结构性信息}
\addcontentsline{toc}{subsection}{5.3 字母层面的结构性信息}


\CodeInline{adjacent\_question\_relation\_stats.csv} 中，\CodeInline{Q40/Q41} 的 \CodeInline{within\_2\_letters\_pct = 3.57\%}，\CodeInline{Q43/Q44} 为 \CodeInline{0\%}，\CodeInline{Q44/Q45} 为 \CodeInline{7.14\%}。这三组都很少在字母序上相距两格以内。结合前面的标准化位置结果，相邻题既没有稳定落在相近段落，字母顺序也经常分离，因此不能沿题号顺序推测原文位置。


\phantomsection
\section*{6. 单题角色重估}
\addcontentsline{toc}{section}{6. 单题角色重估}


\phantomsection
\subsection*{6.1 \CodeInline{Q36}：位置锚，不是起手锚}
\addcontentsline{toc}{subsection}{6.1 Q36：位置锚，不是起手锚}


\CodeInline{question\_order\_profile.csv} 中，\CodeInline{Q36} 的 \CodeInline{median\_rank = 5.0}、\CodeInline{mean\_rank = 4.96}，\CodeInline{earliest\_pct = 0\%}，进入最早两题的比例也只有 \CodeInline{top2\_earliest\_pct = 8.93\%}。它的单题位置相对集中，却从未成为样本中的最早题。因而 \CodeInline{Q36} 更适合在已有阅读线索下帮助定区，而不是固定为每套题的第一步。


\phantomsection
\subsection*{6.2 \CodeInline{Q37}：前段奇数锚}
\addcontentsline{toc}{subsection}{6.2 Q37：前段奇数锚}


\CodeInline{Q37} 的 \CodeInline{earliest\_pct = 10.71\%}，进入最早两题和最早三题的比例分别为 \CodeInline{30.36\%} 与 \CodeInline{44.64\%}。这些比例都高于 \CodeInline{Q36}，所以在前段出现清晰题干线索时，\CodeInline{Q37} 更值得优先核对；这仍是候选顺序的调整，不是无条件规定起手题号。


\phantomsection
\subsection*{6.3 \CodeInline{Q38}：摆动题、反推题}
\addcontentsline{toc}{subsection}{6.3 Q38：摆动题、反推题}


\CodeInline{Q38} 同时有 \CodeInline{earliest\_pct = 19.64\%} 和 \CodeInline{top2\_latest\_pct = 19.64\%}，\CodeInline{median\_rank = 5.5} 也落在中间。这不是一个稳定靠前或稳定靠后的题：仅凭题号去找，搜索成本偏高；可是一旦依据题干把它定位下来，它所在的异常带又会明显改变其他题的候选范围。因此，\CodeInline{Q38} 的价值主要出现在“已经做出以后”，而不是起手阶段。


\phantomsection
\subsection*{6.4 \CodeInline{Q41} 与 \CodeInline{Q45}：后段双锚}
\addcontentsline{toc}{subsection}{6.4 Q41 与 Q45：后段双锚}


\CodeInline{Q41} 的 \CodeInline{top2\_latest\_pct = 33.93\%}，\CodeInline{Q45} 的 \CodeInline{top3\_latest\_pct = 44.64\%}。两道题都偏向后段，但承担的角色不同：前者更像后段的终点锚，后者提供的是较宽的后段方向，不能因为题号最大就默认它一定占据最后一个答案位置。


\phantomsection
\subsection*{6.5 \CodeInline{Q44}：中后段总枢纽}
\addcontentsline{toc}{subsection}{6.5 Q44：中后段总枢纽}


\CodeInline{Q44} 的 \CodeInline{median\_rank = 4.0}，说明它自己并不天然靠后；真正使它成为枢纽的是联动数量：\CodeInline{rule\_count\_gain10 = 24}，\CodeInline{rule\_count\_gain15 = 16}。也就是说，\CodeInline{Q44} 未必用来收尾，却常常在确认后改变多道题的搜索方向。


\phantomsection
\section*{7. 题间无条件相关性}
\addcontentsline{toc}{section}{7. 题间无条件相关性}


\phantomsection
\subsection*{7.1 相邻题通常不是就近题}
\addcontentsline{toc}{subsection}{7.1 相邻题通常不是就近题}


根据 \CodeInline{question\_pair\_relation\_stats.csv} 与 \CodeInline{adjacent\_question\_relation\_stats.csv}：


最典型的四组相邻题都没有表现出“题号挨着，段落也挨着”。\CodeInline{Q40/Q41} 的 \CodeInline{spearman\_pct = -0.2106}，两题落在相差不超过 2 段的比例只有 \CodeInline{3.57\%}，中位段落差为 \CodeInline{4.0}；\CodeInline{Q41/Q42} 的相应三项指标为 \CodeInline{-0.3831}、\CodeInline{8.93\%} 和 \CodeInline{-4.0}。后两组更明显：\CodeInline{Q43/Q44} 的 \CodeInline{within\_2\_paragraph\_pct = 0\%}，\CodeInline{spearman\_pct = -0.4364}，中位段落差为 \CodeInline{-4.0}；\CodeInline{Q44/Q45} 的相关系数降到 \CodeInline{-0.5107}，近邻比例只有 \CodeInline{7.14\%}，中位段落差为 \CodeInline{3.5}。因此，“相邻题顺手在附近找”并没有数据支持。


\phantomsection
\subsection*{7.2 隔一题的同奇偶链更强}
\addcontentsline{toc}{subsection}{7.2 隔一题的同奇偶链更强}


隔一题的组合呈现出不同的形态。\CodeInline{Q37/Q39} 的 \CodeInline{spearman\_pct = 0.2981}，近邻比例为 \CodeInline{51.79\%}，同半区比例为 \CodeInline{60.71\%}；\CodeInline{Q36/Q38} 的三项指标分别为 \CodeInline{0.2863}、\CodeInline{41.07\%} 和 \CodeInline{53.57\%}；\CodeInline{Q43/Q45} 的相关系数虽然只有 \CodeInline{0.2078}，同半区比例却达到 \CodeInline{67.86\%}。三组的中位段落差都是 \CodeInline{1.0}。这些数并不意味着隔一题必然相邻，却足以改变联动顺序：若要借一道题缩小另一道题的范围，应当先检查同奇偶链，而不是只按题号寻找相邻题。


\phantomsection
\subsection*{7.3 字母层面也支持“相邻题分离，隔一题更近”}
\addcontentsline{toc}{subsection}{7.3 字母层面也支持“相邻题分离，隔一题更近”}


\CodeInline{adjacent\_question\_relation\_stats.csv} 和补充复核给出了另一种核对方式：相邻题字母差绝对值 \CodeInline{<= 2} 的比例整体只有 \CodeInline{7.14\%}，隔一题时却达到 \CodeInline{42.19\%}。这与位置百分比的结果相互印证。相邻题不仅在标准化位置上常常分开，换成原文段落字母后也没有出现稳定近邻；隔一题关系虽然仍不能当成硬规则，却更值得在候选收窄时保留。


\phantomsection
\section*{8. 条件联动：这轮最有价值的新信息}
\addcontentsline{toc}{section}{8. 条件联动：这轮最有价值的新信息}


\phantomsection
\subsection*{8.1 跨区间移动规则}
\addcontentsline{toc}{subsection}{8.1 跨区间移动规则}


\phantomsection
\subsubsection*{规则 A：\CodeInline{Q39 = 0-20\% -> Q41 = 60-100\%}}
\addcontentsline{toc}{subsubsection}{规则 A：Q39 = 0-20\% -> Q41 = 60-100\%}


这条规则有 \CodeInline{support\_count = 8} 个条件样本。把 \CodeInline{Q41} 的最佳 40\% 搜索窗放在 \CodeInline{60-100\%} 后，覆盖率由无条件状态提高到 \CodeInline{100.0\%}，\CodeInline{gain\_40\_pct = +55.36}。样本量不大，但窗口移动得很彻底：遇到 \CodeInline{Q39} 已经落入前段的文章，可以降低 \CodeInline{Q41} 前半区候选的优先级，再回到题干判断后段的具体位置，而不是直接把它判成某一段。


\phantomsection
\subsubsection*{规则 B：\CodeInline{Q44 = 0-20\% -> Q45 = 60-100\%}}
\addcontentsline{toc}{subsubsection}{规则 B：Q44 = 0-20\% -> Q45 = 60-100\%}


\CodeInline{Q44} 落在前 20\% 的样本共有 \CodeInline{12} 个，其中 \CodeInline{Q45} 落入 \CodeInline{60-100\%} 的比例为 \CodeInline{91.67\%}，相对无条件窗口提高 \CodeInline{32.74} 个百分点。这里没有出现绝对覆盖，但方向足够一致，实际做题时应当先查后段；若题干在后段找不到对应内容，再保留少数例外，而不是反过来从前段重新遍历。


\phantomsection
\subsection*{8.2 搜索窗口压缩规则}
\addcontentsline{toc}{subsection}{8.2 搜索窗口压缩规则}


\phantomsection
\subsubsection*{规则 C：\CodeInline{Q37 = 0-20\% -> Q39 = 10-50\%}}
\addcontentsline{toc}{subsubsection}{规则 C：Q37 = 0-20\% -> Q39 = 10-50\%}


\CodeInline{Q37} 落在前段的 \CodeInline{17} 个样本中，\CodeInline{Q39} 有 \CodeInline{88.24\%} 落进 \CodeInline{10-50\%}，窗口覆盖率提高 \CodeInline{27.53} 个百分点。无条件分布原本较散，加入 \CodeInline{Q37} 的位置后才收成前中段的一条连续窗口，因此这条规则适合减少阅读范围，但仍不能越过题干把答案定死。


\phantomsection
\subsubsection*{规则 D：\CodeInline{Q42 = 20-40\% -> Q44 = 0-40\%}}
\addcontentsline{toc}{subsubsection}{规则 D：Q42 = 20-40\% -> Q44 = 0-40\%}


当 \CodeInline{Q42} 位于 \CodeInline{20-40\%} 时，\CodeInline{12} 个条件样本中有 \CodeInline{91.67\%} 的 \CodeInline{Q44} 落在前 40\%，覆盖率提高 \CodeInline{39.88} 个百分点。这个变化把 \CodeInline{Q44} 从“可能在多个区域”推回前半区，因而应当先向前查；同一个 \CodeInline{Q42} 位置还会影响 \CodeInline{Q45}，两条信息可以同时使用。


\phantomsection
\subsubsection*{规则 E：\CodeInline{Q42 = 20-40\% -> Q45 = 50-90\%}}
\addcontentsline{toc}{subsubsection}{规则 E：Q42 = 20-40\% -> Q45 = 50-90\%}


还是这 \CodeInline{12} 个样本，\CodeInline{Q45} 落在 \CodeInline{50-90\%} 的比例同样达到 \CodeInline{91.67\%}，\CodeInline{gain\_40\_pct = +32.74}。这时 \CodeInline{Q42} 提供的不是一个孤立答案，而是一组分流信息：\CodeInline{Q44} 往前查，\CodeInline{Q45} 往中后段查。两道题若都挤在同一区域，反而应当重新检查已有匹配。


\phantomsection
\subsection*{8.3 反向排除规则}
\addcontentsline{toc}{subsection}{8.3 反向排除规则}


\phantomsection
\subsubsection*{规则 F：\CodeInline{Q41 = 80-100\% -> Q40 = 20-60\%}}
\addcontentsline{toc}{subsubsection}{规则 F：Q41 = 80-100\% -> Q40 = 20-60\%}


\CodeInline{Q41} 进入最后 20\% 的情况有 \CodeInline{18} 个，\CodeInline{Q40} 落入 \CodeInline{20-60\%} 的覆盖率为 \CodeInline{88.89\%}，比无条件窗口高 \CodeInline{37.10} 个百分点。这里最有用的不是预测一个精确段落，而是排除“相邻题也在末段”的顺手判断：\CodeInline{Q41} 已经靠后时，\CodeInline{Q40} 应先回到中段寻找。


\phantomsection
\subsubsection*{规则 G：\CodeInline{Q43 = 60-80\% -> Q44 = 0-40\%}}
\addcontentsline{toc}{subsubsection}{规则 G：Q43 = 60-80\% -> Q44 = 0-40\%}


\CodeInline{Q43} 落在 \CodeInline{60-80\%} 的 \CodeInline{11} 个样本中，\CodeInline{Q44} 有 \CodeInline{90.91\%} 回到前 40\%，窗口增益为 \CodeInline{39.12}。这不是两题同向移动的链条，而是一个相互拉开的组合：\CodeInline{Q43} 已占住后段后，继续在附近寻找 \CodeInline{Q44} 的收益很低。


\phantomsection
\subsubsection*{规则 H：\CodeInline{Q37 = 60-80\% -> Q36 = 10-50\%}}
\addcontentsline{toc}{subsubsection}{规则 H：Q37 = 60-80\% -> Q36 = 10-50\%}


\CodeInline{Q37} 已在 \CodeInline{60-80\%} 时共有 \CodeInline{15} 个样本，\CodeInline{Q36} 全部落入 \CodeInline{10-50\%}，对应 \CodeInline{gain\_40\_pct = +28.57}。因此可以先把 \CodeInline{Q36} 的候选放在前中段；不过 \CodeInline{10-50\%} 仍覆盖了较宽的区域，这条规则只负责删去明显不合适的后段，不承担独立作答。


\phantomsection
\section*{9. 继续深挖后，新增的高层结构}
\addcontentsline{toc}{section}{9. 继续深挖后，新增的高层结构}


\phantomsection
\subsection*{9.1 边缘带最有信息量}
\addcontentsline{toc}{subsection}{9.1 边缘带最有信息量}


把高增益规则逐条对照后可以发现，最明显的窗口变化几乎都由 \CodeInline{0-20\%} 或 \CodeInline{80-100\%} 触发。中间带中的位置变化大多仍表现为方向性倾向，到了两端，候选才会成片地移向另一侧。因此，边缘带应当优先用来删减搜索范围；中间带更适合给候选排序，不能承担同样强度的排除。


\phantomsection
\subsection*{9.2 很多题对不是线性链条，而是镜像跷跷板}
\addcontentsline{toc}{subsection}{9.2 很多题对不是线性链条，而是镜像跷跷板}


\CodeInline{Q41/Q40} 和 \CodeInline{Q43/Q44} 是两个典型组合。它们不是“前题越前，后题越后”的单向链条，而更接近两端相互让位：\CodeInline{Q41} 靠后时，\CodeInline{Q40} 回到中段；\CodeInline{Q43} 靠后时，\CodeInline{Q44} 转向前半区。把这种关系看成方向相反的条件移动，比笼统地说两题负相关更接近考场中的实际用法。


\phantomsection
\subsection*{9.3 \CodeInline{Q38} 是“难做却好反推”的代表题}
\addcontentsline{toc}{subsection}{9.3 Q38 是“难做却好反推”的代表题}


\CodeInline{Q38} 的无条件分布不稳，但它一旦落在异常带，会同时收窄 \CodeInline{Q37}、\CodeInline{Q39}、\CodeInline{Q40} 和 \CodeInline{Q45} 的候选。这个反差决定了它不适合凭历史位置起手，却适合在前置条件形成后重新处理；定位成功后，应当立即把它的位置传给相关题，而不是只填完一题就继续顺序作答。


\phantomsection
\subsection*{9.4 \CodeInline{Q44} 是中后段总枢纽}
\addcontentsline{toc}{subsection}{9.4 Q44 是中后段总枢纽}


\CodeInline{Q44} 的特点不在某一条孤立规则，而在于它会同时改变 \CodeInline{Q45}、\CodeInline{Q41}、\CodeInline{Q43}、\CodeInline{Q42} 和 \CodeInline{Q39} 的候选顺序。因而判断 \CodeInline{Q44} 的价值时，不能只看它自身落点是否稳定；只要题干证据足够，提前确认它往往能减少后续几道题的重复搜索。


\phantomsection
\section*{10. 图表和统计结果的正确用法}
\addcontentsline{toc}{section}{10. 图表和统计结果的正确用法}


\phantomsection
\subsection*{10.1 正确流程}
\addcontentsline{toc}{subsection}{10.1 正确流程}


\begin{enumerate}
\item 先用单题粗分布确定前、中、后哪一块值得搜索。
\item 把已占用的容量记下来，某一带挤满后就把下一题移开。
\item 题号关系只作参考，不把相邻题当成原文中的邻近题。
\item 搜索具体题目时，才使用对应的条件联动规则。
\end{enumerate}


\phantomsection
\subsection*{10.2 常见误用}
\addcontentsline{toc}{subsection}{10.2 常见误用}


\begin{enumerate}
\item 把箱线图的中位数当成某题的固定答案位置。
\item 把字母分布图误当成“押选项字母”的图。
\item 看到题号相邻，就顺手把原文位置也排成相邻。
\item 只盯着单题热区，忘记检查整套容量是否已经超载。
\end{enumerate}


\phantomsection
\section*{11. 统计局限}
\addcontentsline{toc}{section}{11. 统计局限}


\phantomsection
\subsection*{11.1 小样本条件规则不能无限外推}
\addcontentsline{toc}{subsection}{11.1 小样本条件规则不能无限外推}


有些条件规则只有 \CodeInline{8}、\CodeInline{9} 或 \CodeInline{10} 个样本。这样的结果可以改变候选的检查顺序，却不适合脱离触发条件继续外推，更不能因为历史覆盖率较高就跳过题干核对。


\phantomsection
\subsection*{11.2 双峰题会让均值失真}
\addcontentsline{toc}{subsection}{11.2 双峰题会让均值失真}


\CodeInline{Q37 / Q41 / Q45} 都有双峰或多峰倾向，只看均值或中位数会把两个高发区抹平。


\phantomsection
\subsection*{11.3 图表更适合排除，不适合单独定点}
\addcontentsline{toc}{subsection}{11.3 图表更适合排除，不适合单独定点}


本报告中的图表主要服务于粗筛、压缩窗口和排除冲突候选。它们能告诉考生“先去哪里找”和“哪里暂时不用找”，但不能越过内容匹配直接押中某个精确段落。


\phantomsection
\section*{12. 综合实战策略总结}
\addcontentsline{toc}{section}{12. 综合实战策略总结}


\phantomsection
\subsection*{12.1 总体原则}
\addcontentsline{toc}{subsection}{12.1 总体原则}


做 Section B 时，10 题之间会互相提供位置线索，前面确认的一题会改变后面的搜索范围。可以按下面的顺序处理：


\begin{enumerate}
\item 先找容易定位的锚点题
\item 用题间结构把结果传给别题
\item 随手记住前区、后区和单带的容量
\item 题目落在边缘带时及时调整搜索范围
\end{enumerate}


\phantomsection
\subsection*{12.2 最适合优先关注的题}
\addcontentsline{toc}{subsection}{12.2 最适合优先关注的题}


按单题的稳定性和联动作用，前段先关注 \CodeInline{Q37}、\CodeInline{Q39}；\CodeInline{Q36} 主要帮助判断区域，\CodeInline{Q41}、\CodeInline{Q45} 用来切入后段。若 \CodeInline{Q44} 的题干较容易定位，应当及时处理，因为它确认后能把信息传给多道中后段题。


\phantomsection
\subsection*{12.3 最不适合一上来裸做的题}
\addcontentsline{toc}{subsection}{12.3 最不适合一上来裸做的题}


\CodeInline{Q38}、\CodeInline{Q42} 和 \CodeInline{Q43} 的无条件位置都较分散，不适合一开始只凭历史分布硬找。等锚点落位、条件窗口形成后再回头处理，通常可以少读一部分无效段落。


\phantomsection
\section*{13. 这轮最该背的 14 条动作卡片}
\addcontentsline{toc}{section}{13. 这轮最该背的 14 条动作卡片}


\begin{enumerate}
\item 前 \CodeInline{40\%} 通常只有 \CodeInline{4-5} 题，达到这个数后就该把目光移开。
\item 单个 \CodeInline{20\%} 带通常最多 \CodeInline{3} 题，已有 \CodeInline{3} 题时优先检查其他带。
\item \CodeInline{Q36} 是位置锚，不是默认起手锚。
\item 前段起手先盯 \CodeInline{Q37}，不要把 \CodeInline{Q36} 当成固定起手题。
\item \CodeInline{Q41} 比 \CodeInline{Q45} 更偏向后段终点。
\item \CodeInline{Q44} 不是尾题，它是高杠杆联动题。
\item \CodeInline{Q39} 一旦落在 \CodeInline{0-20\%}，\CodeInline{Q41} 直接优先去最后 \CodeInline{40\%} 找。
\item \CodeInline{Q42} 一旦落在 \CodeInline{20-40\%}，\CodeInline{Q44} 优先回前 \CodeInline{40\%}。
\item \CodeInline{Q41} 落在 \CodeInline{80-100\%} 时，\CodeInline{Q40} 就不要继续在附近找。
\item \CodeInline{Q44} 落在 \CodeInline{0-20\%} 时，\CodeInline{Q45} 优先查 \CodeInline{60-100\%}。
\item 边缘带通常比中间带更能缩小下一道题的范围。
\item 有些题对不是链条，而是跷跷板，尤其 \CodeInline{Q41/Q40}、\CodeInline{Q43/Q44}。
\item \CodeInline{Q38} 不是好起手题，但它是“难做却好反推”的代表题。
\item \CodeInline{Q44} 一旦确认，就马上把信息传给 \CodeInline{Q45 / Q41 / Q43 / Q42}。
\end{enumerate}


\phantomsection
\section*{14. 最终结论}
\addcontentsline{toc}{section}{14. 最终结论}


Section B 的位置规律不能靠一个均值概括。单题基础分布先给出大区，题号角色决定哪些结果值得优先利用，整套容量用于发现局部拥挤；等到锚点落位后，无条件题间结构、条件窗口和反向排除才开始发挥作用。这几类信息不是并排罗列的六条规则，而是前一层为后一层提供条件。

实际执行时，可以把它记成下面这段话：\textbf{先用粗分布确定大区，再用容量校正并确认锚点；出现边缘带后立即收窄窗口。\CodeInline{Q36} 用来定区，\CodeInline{Q44} 用来传递条件，\CodeInline{Q38} 用来反推；前段查 \CodeInline{Q37/Q39}，后段查 \CodeInline{Q41/Q45}，不要按 \CodeInline{36 -> 37 -> ... -> 45} 顺序硬扫。}


\phantomsection
\subsection*{14.1 最终统一实战口径}
\addcontentsline{toc}{subsection}{14.1 最终统一实战口径}


附录中的模型会给出不同指标下的“最优”组合，考场执行时不能把这些最优值混成一个固定题序。起手先在 \CodeInline{Q37 / Q39 / Q41 / Q45} 中选择题干最容易定位的一题，\CodeInline{Q36} 留作定区；\CodeInline{Q38 / Q42 / Q43} 若没有明显文本线索，暂时跳过。

锚点落位后再启用联动。\CodeInline{Q39} 很靠前，就先把 \CodeInline{Q41} 推向后段；\CodeInline{Q42} 在 \CodeInline{20-40\%}，就把 \CodeInline{Q44} 拉回前 40\%；\CodeInline{Q44} 一旦确认，随即检查 \CodeInline{Q45 / Q41 / Q43 / Q42} 的候选是否需要调整。收尾阶段可以降低已用段落的优先级，却不能绝对排除，因为 Section B 明确允许同一段被多次选择。

统计模型在这里始终只作辅助。\CodeInline{Q38+Q40+Q42} 的三锚组合信息互补，\CodeInline{Q41} 的信息增益较大，这些结果说明它们做出来以后有用，并不等于第一题必须从它们开始。


\phantomsection
\section*{15. 附录 A：sample.txt 二次核对汇总}
\addcontentsline{toc}{section}{15. 附录 A：sample.txt 二次核对汇总}


\phantomsection
\subsection*{15.1 结论}
\addcontentsline{toc}{subsection}{15.1 结论}


本轮对 \CodeInline{sample.txt} 纳入答案库的全部 \CodeInline{59} 篇 Section B 文章、共 \CodeInline{590} 题（\CodeInline{Q36-Q45}）完成二次逐题复核。


结果如下：


\begin{itemize}
\item 复核标题数：\CodeInline{59}
\item 复核题目数：\CodeInline{590}
\item 发现需要修正的标题数：\CodeInline{0}
\item 发现需要修正的题目数：\CodeInline{0}
\end{itemize}


也就是说，当前 \CodeInline{MANUAL\_ANSWERS} 中对应 \CodeInline{sample.txt} 的答案，在这一轮逐题复核后，\textbf{没有发现需要改动的题号}。


\phantomsection
\subsection*{15.2 核对范围}
\addcontentsline{toc}{subsection}{15.2 核对范围}


核对材料中有 \CodeInline{56} 篇来自标准 \CodeInline{section\_b\_texts} 题文包，另有 \CodeInline{3} 篇只能沿 \CodeInline{sample.txt} /\allowbreak{} OCR 链路补入，分别是 \CodeInline{The Curious Case of the Tree That Owns Itself}、\CodeInline{Blame your worthless workdays on meeting recovery syndrome} 和 \CodeInline{These are the habits to avoid if you want to make a behavior change}。


\phantomsection
\subsection*{15.3 方法}
\addcontentsline{toc}{subsection}{15.3 方法}


核对没有直接沿用 \CodeInline{sample.txt} 的现成答案，实际做法是：


\begin{enumerate}
\item 每篇文章单独生成逐题复核包。
\item 对照全文段落核对 \CodeInline{36-45} 的题干。
\item 输出每篇文章的答案串，并标出需要修正的题目。
\item 合并各批次结果，检查答案库中是否还有不一致条目。
\end{enumerate}


\phantomsection
\subsection*{15.4 结果文件}
\addcontentsline{toc}{subsection}{15.4 结果文件}


逐批复核结果在：


\begin{itemize}
\item \CodeInline{cet 6/\_analysis\_output/recheck\_reviews/batch1\_review.md}
\item \CodeInline{cet 6/\_analysis\_output/recheck\_reviews/batch2\_review.md}
\item \CodeInline{cet 6/\_analysis\_output/recheck\_reviews/batch3\_review.md}
\item \CodeInline{cet 6/\_analysis\_output/recheck\_reviews/batch4\_review.md}
\item \CodeInline{cet 6/\_analysis\_output/recheck\_reviews/batch5\_review.md}
\item \CodeInline{cet 6/\_analysis\_output/recheck\_reviews/batch6\_review.md}
\end{itemize}


总汇总表在：


\begin{itemize}
\item \CodeInline{cet 6/\_analysis\_output/recheck\_review\_summary.csv}
\end{itemize}


复核包在：


\begin{itemize}
\item \CodeInline{cet 6/\_analysis\_output/recheck\_packets/}
\end{itemize}


\phantomsection
\subsection*{15.5 当前状态}
\addcontentsline{toc}{subsection}{15.5 当前状态}


这轮二次核对没有发现需要修正的答案，因而 \CodeInline{MANUAL\_ANSWERS} 保持不变，也没有触发统计图表重跑。正文引用的统计结论和相应 \CodeInline{md} 报告仍对应当前答案库。


\phantomsection
\section*{16. 附录 B：相关性速记版}
\addcontentsline{toc}{section}{16. 附录 B：相关性速记版}


\phantomsection
\subsection*{16.1 这轮最该记的核心结构}
\addcontentsline{toc}{subsection}{16.1 这轮最该记的核心结构}


这一轮保留了“相邻题弱、隔一题链较强”的判断，但它已不再是策略的全部。对实际作答更有用的是把容量、边缘带、条件窗口和反向排除接起来：\CodeInline{Q36} 负责定区而不是固定起手，\CodeInline{Q44} 负责传递中后段条件，\CodeInline{Q38} 本身难以裸做，却适合在定位后反推；\CodeInline{Q41/Q40} 与 \CodeInline{Q43/Q44} 还表现出相反方向移动的关系。


\phantomsection
\subsection*{16.2 整套容量约束}
\addcontentsline{toc}{subsection}{16.2 整套容量约束}


\CodeInline{question\_capacity\_summary.csv} 显示，前 \CodeInline{40\%} 最常见的是 \CodeInline{4} 题（\CodeInline{39/59 = 66.10\%}），落在 \CodeInline{4-5} 题的共有 \CodeInline{51/59 = 86.44\%}；前 \CodeInline{50\%} 容纳 \CodeInline{5-6} 题的比例为 \CodeInline{52/59 = 88.14\%}。在更小的区间上，单个 \CodeInline{20\%} 带最多装 \CodeInline{3} 题的情况占 \CodeInline{46/59 = 77.97\%}，装到 \CodeInline{4} 题的只有 \CodeInline{3/59 = 5.08\%}。所以前区接近 4--5 题时应当主动转向中后段；某个 20\% 带已经命中 3 题时，继续把新答案塞进去需要更强的题干证据。


\phantomsection
\subsection*{16.3 题号角色重估}
\addcontentsline{toc}{subsection}{16.3 题号角色重估}


\CodeInline{question\_order\_profile.csv} 进一步区分了几道题的角色。\CodeInline{Q36} 的 \CodeInline{median\_rank = 5.0}，且 \CodeInline{earliest\_pct = 0\%}，适合帮助定区，却没有理由固定为第一题。\CodeInline{Q37} 的 \CodeInline{top2\_earliest\_pct = 30.36\%}，是前段更值得留意的奇数锚。\CodeInline{Q38} 的最早与最晚两端比例都是 \CodeInline{19.64\%}，体现的是摆动而非稳定方向。

后段中，\CodeInline{Q41} 的 \CodeInline{top2\_latest\_pct = 33.93\%}，比 \CodeInline{Q45} 更像终点锚；\CodeInline{Q45} 的 \CodeInline{top3\_latest\_pct = 44.64\%}，仍然值得从后段切入，但不能被当作唯一尾题。\CodeInline{Q44} 的 \CodeInline{median\_rank = 4.0}，它自身并不偏尾，价值来自确认后对其他题的带动。


\phantomsection
\subsection*{16.4 杠杆最大的题：Q44}
\addcontentsline{toc}{subsection}{16.4 杠杆最大的题：Q44}


\CodeInline{anchor\_leverage\_summary.csv} 中，\CodeInline{Q44} 的 \CodeInline{rule\_count\_gain10 = 24}、\CodeInline{rule\_count\_gain15 = 16}，\CodeInline{top3\_avg\_gain\_40\_pct = 29.96}。这些指标共同指向同一个用法：\CodeInline{Q44} 不是留到最后再填的题，一旦做出，就应当马上重排 \CodeInline{Q45 / Q41 / Q43 / Q42} 的中后段候选。


\phantomsection
\subsection*{16.5 最值钱的条件联动}
\addcontentsline{toc}{subsection}{16.5 最值钱的条件联动}


来自 \CodeInline{conditional\_window\_gain\_stats.csv}：


\phantomsection
\subsubsection*{强跳带}
\addcontentsline{toc}{subsubsection}{强跳带}


\begin{itemize}
\item \CodeInline{Q39} 在 \CodeInline{0-20\%} 时，\CodeInline{Q41} 落入 \CodeInline{60-100\%} 的覆盖率为 \CodeInline{100\%}
\item \CodeInline{Q44} 在 \CodeInline{0-20\%} 时，\CodeInline{Q45} 落入 \CodeInline{60-100\%} 的覆盖率为 \CodeInline{91.67\%}
\end{itemize}


\phantomsection
\subsubsection*{压缩窗口}
\addcontentsline{toc}{subsubsection}{压缩窗口}


\begin{itemize}
\item \CodeInline{Q37} 在 \CodeInline{0-20\%} 时，\CodeInline{Q39} 被压进 \CodeInline{10-50\%}，覆盖率 \CodeInline{88.24\%}
\item \CodeInline{Q42} 在 \CodeInline{20-40\%} 时，\CodeInline{Q44} 被压进 \CodeInline{0-40\%}，覆盖率 \CodeInline{91.67\%}
\item \CodeInline{Q42} 在 \CodeInline{20-40\%} 时，\CodeInline{Q45} 被压进 \CodeInline{50-90\%}，覆盖率 \CodeInline{91.67\%}
\end{itemize}


\phantomsection
\subsubsection*{反向排除}
\addcontentsline{toc}{subsubsection}{反向排除}


\begin{itemize}
\item \CodeInline{Q41} 在 \CodeInline{80-100\%} 时，\CodeInline{Q40} 被压到 \CodeInline{20-60\%}，覆盖率 \CodeInline{88.89\%}
\item \CodeInline{Q43} 在 \CodeInline{60-80\%} 时，\CodeInline{Q44} 被压回 \CodeInline{0-40\%}，覆盖率 \CodeInline{90.91\%}
\item \CodeInline{Q37} 在 \CodeInline{60-80\%} 时，\CodeInline{Q36} 被压进 \CodeInline{10-50\%}，覆盖率 \CodeInline{100\%}
\end{itemize}


\phantomsection
\subsection*{16.6 高层规律}
\addcontentsline{toc}{subsection}{16.6 高层规律}


\phantomsection
\subsubsection*{边缘带最有信息量}
\addcontentsline{toc}{subsubsection}{边缘带最有信息量}


增益较高的规则大多由 \CodeInline{0-20\%} 和 \CodeInline{80-100\%} 触发；中间带更多提供方向，边缘带才经常造成明显的窗口迁移。两者在策略中的权重不应相同。


\phantomsection
\subsubsection*{镜像跷跷板比线性链更重要}
\addcontentsline{toc}{subsubsection}{镜像跷跷板比线性链更重要}


\CodeInline{Q41} 很前时 \CodeInline{Q40} 被推后，\CodeInline{Q41} 很后时 \CodeInline{Q40} 又被拉回前面；\CodeInline{Q43} 与 \CodeInline{Q44} 也有相似的反向移动。这样的题对不能按线性链理解，更适合在一端确认后排除另一端的同区候选。


\phantomsection
\subsubsection*{\CodeInline{Q38} 是好反推题}
\addcontentsline{toc}{subsubsection}{Q38 是好反推题}


\CodeInline{Q38} 本身分散，不适合裸做；一旦落到异常带，却会明显约束 \CodeInline{Q37 / Q39 / Q40 / Q45}。它是典型的后置反推题。


\phantomsection
\subsubsection*{\CodeInline{Q44} 是全局枢纽}
\addcontentsline{toc}{subsubsection}{Q44 是全局枢纽}


\CodeInline{Q44} 不一定最晚出现，但最能切开后段结构。确认后若不立刻向相关题扩散信息，就浪费了它的主要价值。


\phantomsection
\subsection*{16.7 最简策略版}
\addcontentsline{toc}{subsection}{16.7 最简策略版}


\begin{enumerate}
\item 先记容量账本，不要在同一区间过度堆题。
\item 前段先看 \CodeInline{Q37 / Q39}，不要默认从 \CodeInline{Q36} 起手。
\item 后段用 \CodeInline{Q41 / Q45} 切区，不要只盯 \CodeInline{Q45}。
\item \CodeInline{Q44} 一旦做出，立刻联动 \CodeInline{Q45 / Q41 / Q43 / Q42}。
\item \CodeInline{Q38 / Q42 / Q43} 放到锚点落位后再缩窗处理。
\item 看到边缘带时，要把它当高强度条件，不要只当“稍微偏前/\allowbreak{}偏后”。
\end{enumerate}


\phantomsection
\section*{17. 附录 C：考试实战机械版操作手册}
\addcontentsline{toc}{section}{17. 附录 C：考试实战机械版操作手册}


\phantomsection
\subsection*{17.1 这份手册的目标}
\addcontentsline{toc}{subsection}{17.1 这份手册的目标}


这不是分析报告。


这是一份给考生直接照着执行的机械操作说明。


原则只有两个：


\begin{enumerate}
\item 不追求优雅。
\item 尽量把判断写成 if-else，让考生在考场上少思考、多执行。
\end{enumerate}


适用前提：


\begin{itemize}
\item 你已经看过题干和选项。
\item 你大致读完了文章主旨，知道文章大概在讲什么。
\item 你不是逐字精读，而是准备进入“定位+匹配”的做题模式。
\end{itemize}


\phantomsection
\subsection*{17.2 先记住的固定底层规则}
\addcontentsline{toc}{subsection}{17.2 先记住的固定底层规则}


进入做题前，脑子里先装这 8 条：


\begin{enumerate}
\item 前 \CodeInline{40\%} 通常只会装 \CodeInline{4-5} 题。
\item 前 \CodeInline{50\%} 通常只会装 \CodeInline{5-6} 题。
\item 单个 \CodeInline{20\%} 区间通常最多只装 \CodeInline{3} 题。
\item 相邻题通常不是邻近题，不要默认 \CodeInline{40} 做出来后 \CodeInline{41} 就在附近。
\item 真正更适合联动的是隔一题链。
\item \CodeInline{Q36} 更像定区题，不像默认起手题。
\item \CodeInline{Q44} 是中后段总枢纽，一旦做出要立刻利用。
\item \CodeInline{Q38 / Q42 / Q43} 一般不要一上来裸做。
\end{enumerate}


\phantomsection
\subsection*{17.3 开始做题前：先做 30 秒准备}
\addcontentsline{toc}{subsection}{17.3 开始做题前：先做 30 秒准备}


\phantomsection
\subsubsection*{步骤 1：先看 10 个题干，不急着做}
\addcontentsline{toc}{subsubsection}{步骤 1：先看 10 个题干，不急着做}


你要做的不是立刻选答案，而是给题干分组。


把 10 题分成三类：


\begin{enumerate}
\item 明显是定义/\allowbreak{}总述/\allowbreak{}主旨型
\item 明显有强关键词、专有名词、数字、对比结构的
\item 很抽象、很难一下定位的
\end{enumerate}


\phantomsection
\subsubsection*{步骤 2：先在脑子里标记 5 个优先题}
\addcontentsline{toc}{subsubsection}{步骤 2：先在脑子里标记 5 个优先题}


默认优先顺序：


\begin{enumerate}
\item \CodeInline{Q37}
\item \CodeInline{Q39}
\item \CodeInline{Q41}
\item \CodeInline{Q45}
\item \CodeInline{Q36}
\end{enumerate}


解释：


\begin{itemize}
\item \CodeInline{Q37 / Q39} 适合切前段和奇数链
\item \CodeInline{Q41 / Q45} 适合切后段
\item \CodeInline{Q36} 适合确认整体是前中段启动还是中段启动
\end{itemize}


\phantomsection
\subsubsection*{步骤 3：建立容量账本}
\addcontentsline{toc}{subsubsection}{步骤 3：建立容量账本}


拿到题后，你要在草稿纸上只记三样东西：


\begin{itemize}
\item 前 \CodeInline{40\%} 已经落了几题
\item 前 \CodeInline{50\%} 已经落了几题
\item 每个 \CodeInline{20\%} 带各落了几题
\end{itemize}


不需要精确到段号，只要大致知道：


\begin{itemize}
\item \CodeInline{0-20\%}
\item \CodeInline{20-40\%}
\item \CodeInline{40-60\%}
\item \CodeInline{60-80\%}
\item \CodeInline{80-100\%}
\end{itemize}


每带装了几题


\phantomsection
\subsection*{17.4 第一轮：先找锚点题}
\addcontentsline{toc}{subsection}{17.4 第一轮：先找锚点题}


\phantomsection
\subsubsection*{总原则}
\addcontentsline{toc}{subsubsection}{总原则}


第一轮不要贪多。


目标不是一口气做完 10 题，而是先做出 2 到 4 个锚点题，用来切开全文结构。


\phantomsection
\subsubsection*{if-else 规则 1：先从哪里开始读文章}
\addcontentsline{toc}{subsubsection}{if-else 规则 1：先从哪里开始读文章}


如果你粗读后感觉：


\begin{itemize}
\item 文章前几段就开始给观点或案例
\end{itemize}


那么：


\begin{itemize}
\item 先从前 \CodeInline{30\%} 精读扫描
\item 优先盯 \CodeInline{Q37}、\CodeInline{Q39}
\end{itemize}


否则如果你感觉：


\begin{itemize}
\item 前几段在铺背景，真正信息集中在中后段
\end{itemize}


那么：


\begin{itemize}
\item 先从 \CodeInline{40\%} 以后精读扫描
\item 优先盯 \CodeInline{Q41}、\CodeInline{Q45}
\end{itemize}


否则：


\begin{itemize}
\item 先扫 \CodeInline{20-60\%}
\item 优先盯 \CodeInline{Q36}、\CodeInline{Q39}
\end{itemize}


\phantomsection
\subsubsection*{if-else 规则 2：第一题先做谁}
\addcontentsline{toc}{subsubsection}{if-else 规则 2：第一题先做谁}


如果前段有明显关键词题：


\begin{itemize}
\item 先做 \CodeInline{Q37}
\end{itemize}


否则如果后段有明显关键词题：


\begin{itemize}
\item 先做 \CodeInline{Q41} 或 \CodeInline{Q45}
\end{itemize}


否则：


\begin{itemize}
\item 先做 \CodeInline{Q39}
\end{itemize}


\phantomsection
\subsubsection*{if-else 规则 3：\CodeInline{Q36} 什么时候做}
\addcontentsline{toc}{subsubsection}{if-else 规则 3：Q36 什么时候做}


如果你已经大致知道：


\begin{itemize}
\item 文章前中段开始出题
\end{itemize}


那么：


\begin{itemize}
\item 立刻做 \CodeInline{Q36}
\end{itemize}


否则如果你发现：


\begin{itemize}
\item 前段一直找不到明显对应句
\end{itemize}


那么：


\begin{itemize}
\item 不要硬磨 \CodeInline{Q36}
\item 先去做 \CodeInline{Q39 / Q41 / Q45}
\end{itemize}


\phantomsection
\subsection*{17.5 第二轮：锚点题做出来后怎么机械推进}
\addcontentsline{toc}{subsection}{17.5 第二轮：锚点题做出来后怎么机械推进}


这部分最重要。


\phantomsection
\subsubsection*{17.5.1 如果先做出 \CodeInline{Q37}}
\addcontentsline{toc}{subsubsection}{17.5.1 如果先做出 Q37}


\phantomsection
\paragraph*{if \CodeInline{Q37} 在 \CodeInline{0-20\%}}


那么：


\begin{enumerate}
\item 先做 \CodeInline{Q39}
\item \CodeInline{Q39} 优先在 \CodeInline{10-50\%} 内找
\item 不要一上来去找 \CodeInline{Q45}
\item 如果 \CodeInline{Q39} 也很前，再把 \CodeInline{Q41} 直接往后 \CodeInline{40\%} 怀疑
\end{enumerate}


\phantomsection
\paragraph*{if \CodeInline{Q37} 在 \CodeInline{20-40\%}}


那么：


\begin{enumerate}
\item 先做 \CodeInline{Q39}
\item 先扫 \CodeInline{20-60\%}
\item 然后顺手看 \CodeInline{Q43 / Q44}
\item 不要急着处理 \CodeInline{Q38}
\end{enumerate}


\phantomsection
\paragraph*{if \CodeInline{Q37} 在 \CodeInline{60-80\%}}


那么：


\begin{enumerate}
\item 不要以为 \CodeInline{Q36} 也在附近
\item 直接把 \CodeInline{Q36} 往 \CodeInline{10-50\%} 怀疑
\item 然后做 \CodeInline{Q39}
\item 再看 \CodeInline{Q40}
\end{enumerate}


\phantomsection
\subsubsection*{17.5.2 如果先做出 \CodeInline{Q39}}
\addcontentsline{toc}{subsubsection}{17.5.2 如果先做出 Q39}


\phantomsection
\paragraph*{if \CodeInline{Q39} 在 \CodeInline{0-20\%}}


那么：


\begin{enumerate}
\item \CodeInline{Q41} 直接优先怀疑 \CodeInline{60-100\%}
\item 不要在前中段磨 \CodeInline{Q41}
\item 顺手把 \CodeInline{Q40} 也往 \CodeInline{30-70\%} 怀疑
\item \CodeInline{Q44} 也可以优先怀疑 \CodeInline{20-60\%}
\end{enumerate}


\phantomsection
\paragraph*{if \CodeInline{Q39} 在 \CodeInline{20-40\%}}


那么：


\begin{enumerate}
\item 先做 \CodeInline{Q41}
\item 先扫 \CodeInline{10-50\%} 或 \CodeInline{20-60\%}
\item 然后再决定是否跳后段
\end{enumerate}


\phantomsection
\paragraph*{if \CodeInline{Q39} 在 \CodeInline{40-60\%}}


那么：


\begin{enumerate}
\item \CodeInline{Q41} 不一定更后，优先查 \CodeInline{30-70\%}
\item 不要过早把 \CodeInline{Q41} 直接钉死在最后段
\end{enumerate}


\phantomsection
\subsubsection*{17.5.3 如果先做出 \CodeInline{Q41}}
\addcontentsline{toc}{subsubsection}{17.5.3 如果先做出 Q41}


\phantomsection
\paragraph*{if \CodeInline{Q41} 在 \CodeInline{80-100\%}}


那么：


\begin{enumerate}
\item \CodeInline{Q40} 不要在它附近找
\item \CodeInline{Q40} 优先去 \CodeInline{20-60\%}
\item \CodeInline{Q42} 优先去 \CodeInline{10-50\%}
\item \CodeInline{Q44} 也优先去 \CodeInline{0-40\%}
\end{enumerate}


\phantomsection
\paragraph*{if \CodeInline{Q41} 在 \CodeInline{20-40\%}}


那么：


\begin{enumerate}
\item \CodeInline{Q40} 反而优先去 \CodeInline{60-100\%}
\item \CodeInline{Q42} 优先去 \CodeInline{40-80\%}
\item 不要把 \CodeInline{Q40} 和 \CodeInline{Q41} 当成相邻联动题
\end{enumerate}


\phantomsection
\paragraph*{if \CodeInline{Q41} 在 \CodeInline{40-60\%}}


那么：


\begin{enumerate}
\item \CodeInline{Q42} 优先去 \CodeInline{0-40\%}
\item \CodeInline{Q40} 优先去 \CodeInline{0-40\%}
\item 后段结构还没完全切开，不要急着处理 \CodeInline{Q45}
\end{enumerate}


\phantomsection
\subsubsection*{17.5.4 如果先做出 \CodeInline{Q44}}
\addcontentsline{toc}{subsubsection}{17.5.4 如果先做出 Q44}


这是最关键的分支。


\phantomsection
\paragraph*{if \CodeInline{Q44} 在 \CodeInline{0-20\%}}


那么：


\begin{enumerate}
\item 下一题立刻做 \CodeInline{Q45}
\item \CodeInline{Q45} 直接优先去 \CodeInline{60-100\%}
\item 然后再做 \CodeInline{Q41}
\item \CodeInline{Q43} 优先怀疑 \CodeInline{40-80\%}
\item \CodeInline{Q42} 优先怀疑 \CodeInline{10-50\%}
\end{enumerate}


\phantomsection
\paragraph*{if \CodeInline{Q44} 在 \CodeInline{20-40\%}}


那么：


\begin{enumerate}
\item 下一题先做 \CodeInline{Q45}
\item \CodeInline{Q45} 优先去 \CodeInline{50-90\%}
\item 然后看 \CodeInline{Q41}
\item 再看 \CodeInline{Q43}
\end{enumerate}


\phantomsection
\paragraph*{if \CodeInline{Q44} 在 \CodeInline{80-100\%}}


那么：


\begin{enumerate}
\item 不要默认 \CodeInline{Q45} 也继续在最后段
\item 先把 \CodeInline{Q45} 怀疑到 \CodeInline{30-70\%}
\item \CodeInline{Q41} 优先怀疑 \CodeInline{10-50\%}
\item \CodeInline{Q39} 优先怀疑 \CodeInline{20-50\%}
\end{enumerate}


\phantomsection
\subsubsection*{17.5.5 如果先做出 \CodeInline{Q42}}
\addcontentsline{toc}{subsubsection}{17.5.5 如果先做出 Q42}


\phantomsection
\paragraph*{if \CodeInline{Q42} 在 \CodeInline{20-40\%}}


那么：


\begin{enumerate}
\item 立刻做 \CodeInline{Q44}
\item \CodeInline{Q44} 优先去 \CodeInline{0-40\%}
\item 然后做 \CodeInline{Q45}
\item \CodeInline{Q45} 优先去 \CodeInline{50-90\%}
\end{enumerate}


\phantomsection
\paragraph*{if \CodeInline{Q42} 在 \CodeInline{0-20\%}}


那么：


\begin{enumerate}
\item \CodeInline{Q44} 优先怀疑 \CodeInline{10-50\%}
\item \CodeInline{Q41} 不要立刻钉死在最后段
\end{enumerate}


\phantomsection
\paragraph*{if \CodeInline{Q42} 在 \CodeInline{60-80\%}}


那么：


\begin{enumerate}
\item 不要以为 \CodeInline{Q44} 还在前 \CodeInline{40\%}
\item 将 \CodeInline{Q44} 的范围放宽到 \CodeInline{10-50\%}
\item 随后检查 \CodeInline{Q45}
\end{enumerate}


\phantomsection
\subsubsection*{17.5.6 如果先做出 \CodeInline{Q38}}
\addcontentsline{toc}{subsubsection}{17.5.6 如果先做出 Q38}


默认不要先做它。


但如果你已经做出来了，就马上利用它。


\phantomsection
\paragraph*{if \CodeInline{Q38} 在 \CodeInline{60-80\%}}


那么：


\begin{enumerate}
\item \CodeInline{Q37} 优先去 \CodeInline{10-50\%}
\item \CodeInline{Q39} 优先去 \CodeInline{20-60\%}
\item \CodeInline{Q40} 不要直接去最后段
\end{enumerate}


\phantomsection
\paragraph*{if \CodeInline{Q38} 在 \CodeInline{40-60\%}}


那么：


\begin{enumerate}
\item \CodeInline{Q40} 优先去 \CodeInline{40-80\%}
\item \CodeInline{Q45} 也优先去 \CodeInline{40-80\%}
\end{enumerate}


\phantomsection
\paragraph*{if \CodeInline{Q38} 在 \CodeInline{0-20\%}}


那么：


\begin{enumerate}
\item 不要对别题做太强推断
\item 继续回到 \CodeInline{Q37 / Q39 / Q41} 主线
\end{enumerate}


\phantomsection
\subsection*{17.6 第三轮：什么时候该放弃某题，转做下一题}
\addcontentsline{toc}{subsection}{17.6 第三轮：什么时候该放弃某题，转做下一题}


\phantomsection
\subsubsection*{if-else 规则 4：某题扫了很久还没找到}
\addcontentsline{toc}{subsubsection}{if-else 规则 4：某题扫了很久还没找到}


如果某题你已经在怀疑区间内看了两遍，还是没有明显对应句：


\begin{itemize}
\item 不要继续死磨
\item 先回到锚点题
\item 或者做它的联动题
\end{itemize}


例如：


\begin{itemize}
\item \CodeInline{Q44} 做不出，就先做 \CodeInline{Q42} 或 \CodeInline{Q45}
\item \CodeInline{Q41} 做不出，就先做 \CodeInline{Q39} 或 \CodeInline{Q45}
\item \CodeInline{Q38} 做不出，就先别管它
\end{itemize}


\phantomsection
\subsubsection*{if-else 规则 5：什么时候应该跳段}
\addcontentsline{toc}{subsubsection}{if-else 规则 5：什么时候应该跳段}


如果当前段已经明显塞进：


\begin{itemize}
\item 同一个 \CodeInline{20\%} 带 3 题
\end{itemize}


那么：


\begin{itemize}
\item 立刻降低继续在这里找的优先级
\end{itemize}


如果当前前 \CodeInline{40\%} 已经塞进：


\begin{itemize}
\item 4 题
\end{itemize}


那么：


\begin{itemize}
\item 后续题优先怀疑中后段
\end{itemize}


如果前 \CodeInline{50\%} 已经塞进：


\begin{itemize}
\item 6 题
\end{itemize}


那么：


\begin{itemize}
\item 不要再轻易把剩余题往前塞
\end{itemize}


\phantomsection
\subsection*{17.7 给考生的最短执行版}
\addcontentsline{toc}{subsection}{17.7 给考生的最短执行版}


如果你懒得记上面那么多，只记这套：


\phantomsection
\subsubsection*{起手}
\addcontentsline{toc}{subsubsection}{起手}


\begin{enumerate}
\item 粗读全文大意。
\item 把 \CodeInline{Q37 / Q39 / Q41 / Q45} 作为首批候选，从题干最容易定位的一题开始做。
\item \CodeInline{Q36} 用来定区，起手不固定，不要默认第一题就是它。
\end{enumerate}


\phantomsection
\subsubsection*{中途}
\addcontentsline{toc}{subsubsection}{中途}


\begin{enumerate}
\item 做出 \CodeInline{Q44} 后，立刻联动 \CodeInline{Q45 / Q41 / Q43 / Q42}。
\item 做出 \CodeInline{Q39} 很前，就把 \CodeInline{Q41} 往最后 \CodeInline{40\%} 推。
\item 做出 \CodeInline{Q42} 在 \CodeInline{20-40\%}，就把 \CodeInline{Q44} 拉回前 \CodeInline{40\%}。
\item 做出 \CodeInline{Q41} 在 \CodeInline{80-100\%}，就别在附近找 \CodeInline{Q40}。
\end{enumerate}


\phantomsection
\subsubsection*{收尾}
\addcontentsline{toc}{subsubsection}{收尾}


\begin{enumerate}
\item \CodeInline{Q38 / Q42 / Q43} 放到后面做。
\item 某个区间已经塞满，就赶紧跳区。
\item 相邻题不要默认邻近。
\end{enumerate}


\phantomsection
\subsection*{17.8 最后的机械铁律}
\addcontentsline{toc}{subsection}{17.8 最后的机械铁律}


\begin{enumerate}
\item 不要迷信相邻题。
\item 不要一上来死磕 \CodeInline{Q38 / Q42 / Q43}。
\item \CodeInline{Q36} 定区，\CodeInline{Q37} 起手，\CodeInline{Q41/Q45} 切后段，\CodeInline{Q44} 放大联动。
\item 边缘带一出现，立刻提高警惕。
\item 做题顺序不是 \CodeInline{36 -> 37 -> ... -> 45}，而是“先锚点，后联动，最后收尾”。
\end{enumerate}


\phantomsection
\section*{18. 附录 D：深度分析证据汇总（读者整理版）}
\addcontentsline{toc}{section}{18. 附录 D：深度分析证据汇总（读者整理版）}


以下 9 项分析使用 59 篇文章的完整位置数据计算，作用是为正文结论提供证据，不单独生成新的做题流程。计算脚本：\CodeInline{deep\_analysis.py}，完整结果：\CodeInline{deep\_analysis\_results.md}。


阅读本章时需要先区分三种口径：


\begin{enumerate}
\item \textbf{硬事实}：例如题号顺序与段落顺序不单调、相邻题普遍远离、同段重复合法。
\item \textbf{弱倾向}：例如多数试卷 10 题对应 10 个不同段落、首尾段选中率偏低、某些题对在统计上更常同向。
\item \textbf{实战动作}：必须同时考虑题干可定位性、扫描成本和条件窗口，不能把“信息增益最大”直接等同为“考场第一题先做”。
\end{enumerate}


\phantomsection
\subsection*{18.1 前后半区分裂比}
\addcontentsline{toc}{subsection}{18.1 前后半区分裂比}


\begin{center}
\small
\setlength{\tabcolsep}{4pt}
\begin{longtable}{@{}>{\raggedright\arraybackslash}p{0.307\linewidth}>{\raggedright\arraybackslash}p{0.307\linewidth}>{\raggedright\arraybackslash}p{0.307\linewidth}@{}}
\toprule
\textbf{前/\allowbreak{}后} & \textbf{篇数} & \textbf{占比} \\
\midrule
\endfirsthead
\toprule
\textbf{前/\allowbreak{}后} & \textbf{篇数} & \textbf{占比} \\
\midrule
\endhead
5/\allowbreak{}5 & 34 & 60.7\% \\
6/\allowbreak{}4 & 14 & 25.0\% \\
4/\allowbreak{}6 & 7 & 12.5\% \\
7/\allowbreak{}3 & 1 & 1.8\% \\
\bottomrule
\end{longtable}
\normalsize
\end{center}


\begin{itemize}
\item 5/\allowbreak{}5 是最常见分裂比（60.7\%），非 6/\allowbreak{}4
\item 前半区 ≥5 题概率：87.5\%
\item \textbf{实战}：找到 4 题就急着转向后半区 →\allowbreak{} 约 90\% 试卷会漏掉第 5 题
\end{itemize}


各题前半区概率：Q36(66.1\%) > Q37=Q39(60.7\%) > Q42(58.9\%) > Q44(57.1\%) > Q40(53.6\%) > Q38(50.0\%) > Q41=Q43(37.5\%) > Q45(33.9\%)


\phantomsection
\subsection*{18.2 奇偶交替结构}
\addcontentsline{toc}{subsection}{18.2 奇偶交替结构}


\textbf{间距=1（相邻）}：全部 9/\allowbreak{}9 对负相关（ρ = -0.21 \textasciitilde{} -0.51） ✓


\textbf{间距=2（隔一）}：5/\allowbreak{}8 对正相关，3 对为负/\allowbreak{}近零 — \textbf{非完美交替}

\begin{itemize}
\item 正相关（1）：Q36/\allowbreak{}Q38(+0.29),\allowbreak{} Q37/\allowbreak{}Q39(+0.30),\allowbreak{} Q41/\allowbreak{}Q43(+0.12)
\item 正相关（2）：Q42/\allowbreak{}Q44(+0.13),\allowbreak{} Q43/\allowbreak{}Q45(+0.21)
\item 负/\allowbreak{}近零：Q38/\allowbreak{}Q40(-0.03),\allowbreak{} Q39/\allowbreak{}Q41(-0.20),\allowbreak{} Q40/\allowbreak{}Q42(-0.01)
\end{itemize}


\textbf{修正结论}：相邻题跷跷板是铁律（9/\allowbreak{}9），但隔一题同向只是多数倾向（5/\allowbreak{}8），不是结构性保证。中段（Q38-Q42）交替结构最弱。


\phantomsection
\subsection*{18.3 文章段落数分组}
\addcontentsline{toc}{subsection}{18.3 文章段落数分组}


\begin{center}
\small
\setlength{\tabcolsep}{4pt}
\begin{longtable}{@{}>{\raggedright\arraybackslash}p{0.230\linewidth}>{\raggedright\arraybackslash}p{0.230\linewidth}>{\raggedright\arraybackslash}p{0.230\linewidth}>{\raggedright\arraybackslash}p{0.230\linewidth}@{}}
\toprule
\textbf{段落数} & \textbf{篇数} & \textbf{占比} & \textbf{干扰段} \\
\midrule
\endfirsthead
\toprule
\textbf{段落数} & \textbf{篇数} & \textbf{占比} & \textbf{干扰段} \\
\midrule
\endhead
10 & 1 & 1.8\% & 0 \\
11 & 8 & 14.3\% & 1 \\
12 & 7 & 12.5\% & 2 \\
13 & 7 & 12.5\% & 3 \\
14 & 7 & 12.5\% & 4 \\
15 & 10 & 17.9\% & 5 \\
16 & 10 & 17.9\% & 6 \\
17 & 4 & 7.1\% & 7 \\
18 & 2 & 3.6\% & 8 \\
\bottomrule
\end{longtable}
\normalsize
\end{center}


\begin{itemize}
\item 中位/\allowbreak{}均值均为 14 段，干扰段 0\textasciitilde{}8 个
\item 短文(≤12)与长文(≥16)位置分布差异显著：Q38 短文中位 59.1\% vs 长文 23.4\%（差 35.7\%），Q37 短文 21.8\% vs 长文 46.9\%（差 -25.1\%）
\item \textbf{实战}：段落数多的文章，干扰段更多，需更强排除法
\end{itemize}


\phantomsection
\subsection*{18.4 时间窗口稳定性}
\addcontentsline{toc}{subsection}{18.4 时间窗口稳定性}


\begin{center}
\small
\setlength{\tabcolsep}{4pt}
\begin{longtable}{@{}>{\raggedright\arraybackslash}p{0.460\linewidth}>{\raggedright\arraybackslash}p{0.460\linewidth}@{}}
\toprule
\textbf{时代} & \textbf{篇数} \\
\midrule
\endfirsthead
\toprule
\textbf{时代} & \textbf{篇数} \\
\midrule
\endhead
2016-2018 & 18 \\
2019-2021 & 17 \\
2022-2025 & 21 \\
\bottomrule
\end{longtable}
\normalsize
\end{center}


\begin{itemize}
\item 平均跨时代位置漂移：17.5\%
\item 最大漂移：Q45（36.9\%，从 33.1\% 漂到 70.0\%）
\item 漂移 ≤10\% 的题仅 3/\allowbreak{}10（Q36,\allowbreak{} Q39,\allowbreak{} Q42）
\item \textbf{结论}：整体结构基本稳定，但个别题（Q45,\allowbreak{} Q44,\allowbreak{} Q37）有明显时代漂移
\end{itemize}


\phantomsection
\subsection*{18.5 同段重复率：不是错误，也不是硬约束}
\addcontentsline{toc}{subsection}{18.5 同段重复率：不是错误，也不是硬约束}


\begin{itemize}
\item 总题对 2520，同答案字母碰撞仅 6 次（0.24\%）
\item 10 题全部不同字母：52/\allowbreak{}59 = 88.1\%
\item 6 篇存在重复字母，说明“同一段对应多题”在样本中少见但真实存在
\item 题目说明明确允许 \CodeInline{You may choose a paragraph more than once}，因此重复字母\textbf{不是数据错误}，也不是必须回溯修正的证据
\item \textbf{实战}：已用过的段落可以降低优先级，但不能直接排除；只有题干语义明显不符时才能排除
\end{itemize}


\phantomsection
\subsection*{18.6 答案序列单调性}
\addcontentsline{toc}{subsection}{18.6 答案序列单调性}


\begin{itemize}
\item 完全单调递增：0/\allowbreak{}59（0\%）— 没有任何一篇是 Q36→\allowbreak{}Q45 从前到后排列
\item 平均逆序对 20.3/\allowbreak{}45（随机期望 22.5），实际/\allowbreak{}随机比 0.90
\item Kendall τ 均值仍接近 0，τ > 0 篇数 41/\allowbreak{}59
\item \textbf{结论}：题号与段落位置基本无单调关系。不能按题号顺序扫文章
\end{itemize}


\phantomsection
\subsection*{18.7 答案字母分布与弱排除}
\addcontentsline{toc}{subsection}{18.7 答案字母分布与弱排除}


\begin{itemize}
\item 共 18 个不同字母被使用（A\textasciitilde{}R）
\item 高频：G(8.6\%) > D=E(8.4\%) > C(8.0\%) > H(7.9\%)
\item 低频：R(0.2\%),\allowbreak{} Q(0.5\%),\allowbreak{} P(1.1\%)
\item \textbf{策略}：字母分布可用于后期弱排除和警惕干扰段，但不能把已用字母当作绝对不可再选
\end{itemize}


\phantomsection
\subsection*{18.8 干扰段落分析}
\addcontentsline{toc}{subsection}{18.8 干扰段落分析}


\begin{itemize}
\item 59 篇共 236 个干扰段落，每篇平均 4.0 个
\item 位置分布：80-100\%(27.7\%) > 60-80\%(22.9\%) > 0-20\%(21.2\%) > 40-60\%(14.3\%) > 20-40\%(13.9\%)
\item 干扰段中位位置 60.7\%
\item \textbf{实战}：文章后半段（尤其 80-100\%）更可能出现干扰段落 →\allowbreak{} 后段匹配需更谨慎
\end{itemize}


\phantomsection
\subsection*{18.9 多锚点复合条件}
\addcontentsline{toc}{subsection}{18.9 多锚点复合条件}


双锚点高确信规则（样本≥4，目标缩至≤2个位置带）：


\begin{center}
\small
\setlength{\tabcolsep}{4pt}
\begin{longtable}{@{}>{\raggedright\arraybackslash}p{0.172\linewidth}>{\raggedright\arraybackslash}p{0.172\linewidth}>{\raggedright\arraybackslash}p{0.172\linewidth}>{\raggedright\arraybackslash}p{0.172\linewidth}>{\raggedright\arraybackslash}p{0.172\linewidth}@{}}
\toprule
\textbf{锚点条件} & \textbf{目标题} & \textbf{样本} & \textbf{最可能带} & \textbf{命中率} \\
\midrule
\endfirsthead
\toprule
\textbf{锚点条件} & \textbf{目标题} & \textbf{样本} & \textbf{最可能带} & \textbf{命中率} \\
\midrule
\endhead
Q36∈40-60\% ∧ Q44∈80-100\% & Q37 & 4 & 0-20\% & 100\% \\
Q36∈40-60\% ∧ Q44∈80-100\% & Q39 & 4 & 20-40\% & 100\% \\
Q37∈0-20\% ∧ Q44∈0-20\% & Q36 & 4 & 40-60\% & 100\% \\
Q37∈0-20\% ∧ Q44∈0-20\% & Q39 & 4 & 20-40\% & 100\% \\
Q36∈40-60\% ∧ Q41∈20-40\% & Q37 & 4 & 0-20\% & 100\% \\
Q36∈40-60\% ∧ Q37∈0-20\% & Q39 & 11 & 20-40\% & 91\% \\
Q37∈0-20\% ∧ Q43∈40-60\% & Q39 & 7 & 20-40\% & 86\% \\
Q36∈40-60\% ∧ Q44∈0-20\% & Q39 & 6 & 20-40\% & 83\% \\
Q37∈0-20\% ∧ Q44∈20-40\% & Q41 & 6 & 80-100\% & 83\% \\
\bottomrule
\end{longtable}
\normalsize
\end{center}


Q36+Q44 双锚定位效果：18 种组合中，Q39 有 4 种可缩至≤2带（22\%），Q37 和 Q45 各 3 种（17\%）。


\textbf{策略}：先确认 Q36 和 Q44 两个锚点，再用双锚条件定位其余题。


\phantomsection
\section*{19. 附录 E：标准化分析证据（不是新策略版本）}
\addcontentsline{toc}{section}{19. 附录 E：标准化分析证据（不是新策略版本）}


以下分析消除了段落数差异的影响，对附录 D 中的粗糙结论进行修正和深化。这里的“最优”均指某个统计指标下的最优，不等于考场执行顺序。完整计算：\CodeInline{deep\_analysis\_round2.py} →\allowbreak{} \CodeInline{deep\_analysis\_round2.md}。


\phantomsection
\subsection*{19.1 标准化段落位置偏好}
\addcontentsline{toc}{subsection}{19.1 标准化段落位置偏好}


将每段转为相对位置（第 k/\allowbreak{}N），消除"P/\allowbreak{}Q/\allowbreak{}R 少只是因为文章短"的假象：


\begin{center}
\small
\setlength{\tabcolsep}{4pt}
\begin{longtable}{@{}>{\raggedright\arraybackslash}p{0.307\linewidth}>{\raggedright\arraybackslash}p{0.307\linewidth}>{\raggedright\arraybackslash}p{0.307\linewidth}@{}}
\toprule
\textbf{位置桶} & \textbf{选中率} & \textbf{偏差(vs期望70.6\%)} \\
\midrule
\endfirsthead
\toprule
\textbf{位置桶} & \textbf{选中率} & \textbf{偏差(vs期望70.6\%)} \\
\midrule
\endhead
0-10\% & 58.3\% & \textbf{-12.2\%} \\
10-20\% & 76.8\% & +6.3\% \\
20-30\% & 83.1\% & +12.5\% \\
40-50\% & \textbf{83.8\%} & \textbf{+13.3\%} \\
80-90\% & 58.3\% & \textbf{-12.2\%} \\
90-100\% & 58.5\% & \textbf{-12.0\%} \\
\bottomrule
\end{longtable}
\normalsize
\end{center}


\textbf{结论}：出题人显著偏好文章中段（20-50\%）的段落作为答案。首尾段（0-10\%、80-100\%）被选为答案的概率比期望低 12\%，更可能是干扰项。


\phantomsection
\subsection*{19.2 题号排名分析}
\addcontentsline{toc}{subsection}{19.2 题号排名分析}


将 10 个答案按段落位置排序，每题排第几？


\begin{center}
\small
\setlength{\tabcolsep}{4pt}
\begin{longtable}{@{}>{\raggedright\arraybackslash}p{0.172\linewidth}>{\raggedright\arraybackslash}p{0.172\linewidth}>{\raggedright\arraybackslash}p{0.172\linewidth}>{\raggedright\arraybackslash}p{0.172\linewidth}>{\raggedright\arraybackslash}p{0.172\linewidth}@{}}
\toprule
\textbf{题号} & \textbf{中位排名} & \textbf{σ} & \textbf{排第1} & \textbf{排第10} \\
\midrule
\endfirsthead
\toprule
\textbf{题号} & \textbf{中位排名} & \textbf{σ} & \textbf{排第1} & \textbf{排第10} \\
\midrule
\endhead
Q36 & 5.0 & \textbf{1.9} & 0 & 0 \\
Q38 & 5.5 & \textbf{3.2} & 11 & 5 \\
Q41 & 7.0 & 3.0 & 5 & 12 \\
Q45 & 7.0 & 2.6 & 1 & 5 \\
\bottomrule
\end{longtable}
\normalsize
\end{center}


\begin{itemize}
\item Q36 排名最稳定（σ=1.9，永远在中间附近）— 名副其实的"位置锚"
\item Q38 最不稳定（σ=3.2，可排第1也可排第10）— 高风险高回报
\item Q41/\allowbreak{}Q45 中位排第7 — 后段偏好确认
\end{itemize}


\textbf{相邻题排名差}：平均仍在 4-5 个位置附近。Q43/\allowbreak{}Q44 排名差≤2 次数 = \textbf{1/\allowbreak{}59}，依然几乎可以视为绝对互斥。相邻题在文章中几乎从不相邻。


\phantomsection
\subsection*{19.3 段落覆盖结构}
\addcontentsline{toc}{subsection}{19.3 段落覆盖结构}


\begin{itemize}
\item 71.6\% 答案段落与前一个答案段连续（间隔=0）
\item 平均最长连续答案段：5.6 个
\item 覆盖跨度：平均 93.3\%，最小 76.5\%
\end{itemize}


\textbf{含义}：答案段落倾向于连片分布，但覆盖几乎全文。不会出现文章某大段完全没有答案的情况。


\phantomsection
\subsection*{19.4 位置重标定}
\addcontentsline{toc}{subsection}{19.4 位置重标定}


用 letter\_index/\allowbreak{}(pcount-1) 替代原始 pct 后，各题中位数变化 < 1.1\% →\allowbreak{} \textbf{原始 pct 指标足够可靠}，段落数差异不影响核心结论。


\phantomsection
\subsection*{19.5 三题锚组合的信息互补性}
\addcontentsline{toc}{subsection}{19.5 三题锚组合的信息互补性}


\begin{center}
\small
\setlength{\tabcolsep}{4pt}
\begin{longtable}{@{}>{\raggedright\arraybackslash}p{0.230\linewidth}>{\raggedright\arraybackslash}p{0.230\linewidth}>{\raggedright\arraybackslash}p{0.230\linewidth}>{\raggedright\arraybackslash}p{0.230\linewidth}@{}}
\toprule
\textbf{排名} & \textbf{锚组合} & \textbf{残余熵} & \textbf{含义} \\
\midrule
\endfirsthead
\toprule
\textbf{排名} & \textbf{锚组合} & \textbf{残余熵} & \textbf{含义} \\
\midrule
\endhead
\textbf{1} & \textbf{Q38+Q40+Q42} & \textbf{0.183} & 做完后剩余7题每题仅需搜 1.1 个带 \\
2 & Q38+Q43+Q45 & 0.202 & \  \\
3 & Q38+Q40+Q44 & 0.208 & \  \\
\textbf{120} & \textbf{Q36+Q37+Q38} & \textbf{0.679} & 最差 — 前三题全做信息量最低 \\
\bottomrule
\end{longtable}
\normalsize
\end{center}


\textbf{解释}：在“做完三题后剩余题残余熵最低”这个指标下，\textbf{Q38+Q40+Q42} 最优，因为它们分布在不同链条和中段区域，位置互补。这个结论不能直接翻译成“考场第一题先做 Q38”，因为起手还必须考虑题干是否好定位。


\phantomsection
\subsection*{19.6 难度代理排名}
\addcontentsline{toc}{subsection}{19.6 难度代理排名}


综合位置离散度、信息熵、干扰段密度：


\begin{center}
\small
\setlength{\tabcolsep}{4pt}
\begin{longtable}{@{}>{\raggedright\arraybackslash}p{0.920\linewidth}@{}}
\toprule
\textbf{最难 →\allowbreak{} 最易} \\
\midrule
\endfirsthead
\toprule
\textbf{最难 →\allowbreak{} 最易} \\
\midrule
\endhead
Q41 > Q38 > Q42 > Q45 > Q37 > Q44 > Q40 > Q43 > Q39 > Q36 \\
\bottomrule
\end{longtable}
\normalsize
\end{center}


Q36 最易（σ 最小、熵最低） →\allowbreak{} 适合第一个做。Q41 最难 →\allowbreak{} 放最后。


\phantomsection
\subsection*{19.7 每题最佳扫描窗口}
\addcontentsline{toc}{subsection}{19.7 每题最佳扫描窗口}


\begin{center}
\small
\setlength{\tabcolsep}{4pt}
\begin{longtable}{@{}>{\raggedright\arraybackslash}p{0.307\linewidth}>{\raggedright\arraybackslash}p{0.307\linewidth}>{\raggedright\arraybackslash}p{0.307\linewidth}@{}}
\toprule
\textbf{题号} & \textbf{最佳40\%窗口} & \textbf{覆盖率} \\
\midrule
\endfirsthead
\toprule
\textbf{题号} & \textbf{最佳40\%窗口} & \textbf{覆盖率} \\
\midrule
\endhead
Q36 & 15-55\% & 76.8\% \\
Q39 & 20-60\% & 60.7\% \\
Q45 & 52-92\% & 60.7\% \\
Q37 & 11-51\% & 58.9\% \\
Q41 & 50-90\% & 48.2\% \\
Q38 & 0-40\% & 44.6\% \\
\bottomrule
\end{longtable}
\normalsize
\end{center}


Q36 在 15-55\% 覆盖 76.8\% →\allowbreak{} 最值得优先扫描的题。Q38 即使最佳窗口也只覆盖 44.6\% →\allowbreak{} 最难定位。


\phantomsection
\subsection*{19.8 同段重复清单}
\addcontentsline{toc}{subsection}{19.8 同段重复清单}


6 篇存在同字母重复。由于 Section B 题目说明允许同一段被多次选择，这些条目不能仅凭重复就判定为错误；它们更适合作为“重复段落现象”的样本清单：


\begin{center}
\small
\setlength{\tabcolsep}{4pt}
\begin{longtable}{@{}>{\raggedright\arraybackslash}p{0.460\linewidth}>{\raggedright\arraybackslash}p{0.460\linewidth}@{}}
\toprule
\textbf{文章} & \textbf{同段重复} \\
\midrule
\endfirsthead
\toprule
\textbf{文章} & \textbf{同段重复} \\
\midrule
\endhead
The American Workplace Is Broken & B→\allowbreak{}Q37,\allowbreak{}Q45 \\
Who's Really Addicting You to Technology & L→\allowbreak{}Q37,\allowbreak{}Q44 \\
Rich Children and Poor Ones Are Raised V... & G→\allowbreak{}Q36,\allowbreak{}Q40 \\
Slow Hope & I→\allowbreak{}Q37,\allowbreak{}Q44 \\
This man is running 7 marathons & I→\allowbreak{}Q37,\allowbreak{}Q43 \\
The problem with being perfect & E→\allowbreak{}Q36,\allowbreak{}Q43 \\
\bottomrule
\end{longtable}
\normalsize
\end{center}


\phantomsection
\subsection*{19.9 对正文实战策略的补充}
\addcontentsline{toc}{subsection}{19.9 对正文实战策略的补充}


附录 D/\allowbreak{}E 的计算结果还给出了一条“信息互补”原则：


\begin{enumerate}
\item \CodeInline{Q36} 适合做稳定定位锚，尤其在 \CodeInline{15-55\%} 内覆盖率高。
\item \CodeInline{Q38 / Q40 / Q42} 做出来后信息互补性强，适合作为中段锚链，但不应在题干难定位时硬做。
\item \CodeInline{Q41} 难度代理最高，适合在候选减少后处理，而不是只凭信息增益直接起手。
\item 后续用锚链、条件窗口和弱排除处理剩余题。
\end{enumerate}


\textbf{补充原则}：

\begin{itemize}
\item 不只依赖 \CodeInline{Q36+Q44} 双锚，也要把 \CodeInline{Q38/Q40/Q42} 这类中段互补题纳入条件判断。
\item 不按 \CodeInline{Q36→Q37→Q38} 机械连做，因为前三题信息重叠较高。
\item 文章首尾段更可能是干扰项，选中率比期望低约 12\%。
\end{itemize}


\phantomsection
\section*{20. 附录 F：决策模型证据与同段重复校正}
\addcontentsline{toc}{section}{20. 附录 F：决策模型证据与同段重复校正}


本章保留决策树、信息增益和条件窗口证据，用来解释正文策略为什么这样安排。完整计算：\CodeInline{deep\_analysis\_round3.py} →\allowbreak{} \CodeInline{deep\_analysis\_round3.md}。


\phantomsection
\subsection*{20.1 同段重复清单（不是数据错误清单）}
\addcontentsline{toc}{subsection}{20.1 同段重复清单（不是数据错误清单）}


\begin{center}
\small
\setlength{\tabcolsep}{4pt}
\begin{longtable}{@{}>{\raggedright\arraybackslash}p{0.230\linewidth}>{\raggedright\arraybackslash}p{0.230\linewidth}>{\raggedright\arraybackslash}p{0.230\linewidth}>{\raggedright\arraybackslash}p{0.230\linewidth}@{}}
\toprule
\textbf{文章} & \textbf{来源} & \textbf{重复} & \textbf{缺失} \\
\midrule
\endfirsthead
\toprule
\textbf{文章} & \textbf{来源} & \textbf{重复} & \textbf{缺失} \\
\midrule
\endhead
The American Workplace Is Broken & 2016.12第3套 & B→\allowbreak{}Q37,\allowbreak{}Q45 & A,\allowbreak{}G,\allowbreak{}I,\allowbreak{}L \\
Who's Really Addicting You to Technology & 2017.12第1套 & L→\allowbreak{}Q37,\allowbreak{}Q44 & A,\allowbreak{}C,\allowbreak{}H,\allowbreak{}K,\allowbreak{}M,\allowbreak{}N \\
Rich Children and Poor Ones & 2017.6第3套 & G→\allowbreak{}Q36,\allowbreak{}Q40 & A,\allowbreak{}C,\allowbreak{}E,\allowbreak{}I,\allowbreak{}J,\allowbreak{}L,\allowbreak{}N,\allowbreak{}Q \\
Slow Hope & 2020.12第2套 & I→\allowbreak{}Q37,\allowbreak{}Q44 & H,\allowbreak{}J \\
7 marathons on 7 continents & 2022.12第1套 & I→\allowbreak{}Q37,\allowbreak{}Q43 & B,\allowbreak{}J,\allowbreak{}L \\
The problem with being perfect & 2023.6第2套 & E→\allowbreak{}Q36,\allowbreak{}Q43 & D,\allowbreak{}J \\
\bottomrule
\end{longtable}
\normalsize
\end{center}


\textbf{重要修正}：这些重复不能直接判定为录入错误，因为题目说明允许同一段被多次选择。它们真正提醒我们的是：所有基于“一一映射”的策略都必须降级为弱排除，不能作为硬规则。


\phantomsection
\subsection*{20.2 贪心信息增益序列}
\addcontentsline{toc}{subsection}{20.2 贪心信息增益序列}


\begin{center}
\small
\setlength{\tabcolsep}{4pt}
\begin{longtable}{@{}>{\raggedright\arraybackslash}p{0.230\linewidth}>{\raggedright\arraybackslash}p{0.230\linewidth}>{\raggedright\arraybackslash}p{0.230\linewidth}>{\raggedright\arraybackslash}p{0.230\linewidth}@{}}
\toprule
\textbf{步骤} & \textbf{做题} & \textbf{信息增益} & \textbf{角色} \\
\midrule
\endfirsthead
\toprule
\textbf{步骤} & \textbf{做题} & \textbf{信息增益} & \textbf{角色} \\
\midrule
\endhead
1 & \textbf{Q41} & 3.390 & 信息增益最大，但不等于考场起手最优 \\
2 & Q37 & — & 高信息量 \\
3 & Q44 & — & 高信息量 \\
4 & Q43 & — & \  \\
5 & Q39 & — & \  \\
6-10 & Q36 →\allowbreak{} Q45 →\allowbreak{} Q42 →\allowbreak{} Q40 →\allowbreak{} Q38 & — & 排除法完成剩余题 \\
\bottomrule
\end{longtable}
\normalsize
\end{center}


\textbf{统计贪心序列：Q41 →\allowbreak{} Q37 →\allowbreak{} Q44 →\allowbreak{} Q43 →\allowbreak{} Q39 →\allowbreak{} Q36 →\allowbreak{} Q45 →\allowbreak{} Q42 →\allowbreak{} Q40 →\allowbreak{} Q38}


与 Round 2 的锚链分析(Q38+Q40+Q42)角度不同：

\begin{itemize}
\item 锚链分析优化的是"三题组合后的残余熵"
\item 贝叶斯贪心优化的是"每步选择的边际信息增益"
\item 两者并不矛盾，但都不能机械等同于考场起手顺序。Q41 信息增益大，是因为它一旦被确认会强力约束后续题；可它本身也是高难题，裸做成本高
\end{itemize}


\phantomsection
\subsection*{20.3 策略搜索效率量化对比}
\addcontentsline{toc}{subsection}{20.3 策略搜索效率量化对比}


做完 k 题后剩余题平均候选带数：


\begin{center}
\small
\setlength{\tabcolsep}{4pt}
\begin{longtable}{@{}>{\raggedright\arraybackslash}p{0.307\linewidth}>{\raggedright\arraybackslash}p{0.307\linewidth}>{\raggedright\arraybackslash}p{0.307\linewidth}@{}}
\toprule
\textbf{步骤} & \textbf{顺序(36→\allowbreak{}45)} & \textbf{锚链(36→\allowbreak{}38→\allowbreak{}40→\allowbreak{}42→\allowbreak{}44)} \\
\midrule
\endfirsthead
\toprule
\textbf{步骤} & \textbf{顺序(36→\allowbreak{}45)} & \textbf{锚链(36→\allowbreak{}38→\allowbreak{}40→\allowbreak{}42→\allowbreak{}44)} \\
\midrule
\endhead
1 & 4.50 & 4.50 \\
2 & 3.00 & \textbf{2.64} \\
3 & 1.82 & \textbf{1.34} \\
4 & 1.30 & \textbf{1.11} \\
5 & 1.13 & \textbf{1.07} \\
\bottomrule
\end{longtable}
\normalsize
\end{center}


锚链策略到第 3 步就缩至 1.34 带，顺序策略需到第 4 步才达到 1.30。锚链快一步。


\phantomsection
\subsection*{20.4 相邻题位置差的真实分布}
\addcontentsline{toc}{subsection}{20.4 相邻题位置差的真实分布}


\begin{center}
\small
\setlength{\tabcolsep}{4pt}
\begin{longtable}{@{}>{\raggedright\arraybackslash}p{0.230\linewidth}>{\raggedright\arraybackslash}p{0.230\linewidth}>{\raggedright\arraybackslash}p{0.230\linewidth}>{\raggedright\arraybackslash}p{0.230\linewidth}@{}}
\toprule
\textbf{题对} & \textbf{中位差} & \textbf{差>40\% 占比} & \textbf{差<20\% 占比} \\
\midrule
\endfirsthead
\toprule
\textbf{题对} & \textbf{中位差} & \textbf{差>40\% 占比} & \textbf{差<20\% 占比} \\
\midrule
\endhead
Q43/\allowbreak{}Q44 & \textbf{45.8\%} & \textbf{61\%} & \textbf{0\%} \\
Q37/\allowbreak{}Q38 & 44.1\% & 61\% & 14\% \\
Q40/\allowbreak{}Q41 & 44.2\% & 57\% & 4\% \\
Q36/\allowbreak{}Q37 & 34.5\% & 32\% & 9\% \\
\bottomrule
\end{longtable}
\normalsize
\end{center}


\textbf{Q43/\allowbreak{}Q44 是最极端的互斥对}：中位差 45.8\%，0\% 概率在 20\% 以内。这两题在文章中永远隔得很远。


\phantomsection
\subsection*{20.5 段落类型偏好}
\addcontentsline{toc}{subsection}{20.5 段落类型偏好}


\begin{center}
\small
\setlength{\tabcolsep}{4pt}
\begin{longtable}{@{}>{\raggedright\arraybackslash}p{0.307\linewidth}>{\raggedright\arraybackslash}p{0.307\linewidth}>{\raggedright\arraybackslash}p{0.307\linewidth}@{}}
\toprule
\textbf{类型} & \textbf{选率} & \textbf{vs 期望} \\
\midrule
\endfirsthead
\toprule
\textbf{类型} & \textbf{选率} & \textbf{vs 期望} \\
\midrule
\endhead
首段(A) & 58.9\% & \textbf{-12.4\%} \\
第二段(B) & 73.2\% & +1.9\% \\
倒数第二段 & 58.9\% & \textbf{-12.4\%} \\
末段 & 60.7\% & \textbf{-10.6\%} \\
\bottomrule
\end{longtable}
\normalsize
\end{center}


\textbf{首段和末段被选为答案的概率显著低于期望} →\allowbreak{} 文章首段（通常是引言）和末段（通常是总结）更可能是干扰项。


\phantomsection
\subsection*{20.6 条件搜索窗口精选（最窄规则）}
\addcontentsline{toc}{subsection}{20.6 条件搜索窗口精选（最窄规则）}


已知锚点位置带 →\allowbreak{} 目标题 80\% 概率落在的范围：


\begin{center}
\small
\setlength{\tabcolsep}{4pt}
\begin{longtable}{@{}>{\raggedright\arraybackslash}p{0.230\linewidth}>{\raggedright\arraybackslash}p{0.230\linewidth}>{\raggedright\arraybackslash}p{0.230\linewidth}>{\raggedright\arraybackslash}p{0.230\linewidth}@{}}
\toprule
\textbf{条件} & \textbf{目标} & \textbf{窗口} & \textbf{宽度} \\
\midrule
\endfirsthead
\toprule
\textbf{条件} & \textbf{目标} & \textbf{窗口} & \textbf{宽度} \\
\midrule
\endhead
Q44∈60-80\% & Q36 & 28-41\% & \textbf{13\%} \\
Q36∈0-20\% & Q38 & 4-25\% & \textbf{21\%} \\
Q40∈80-100\% & Q39 & 27-53\% & \textbf{26\%} \\
Q38∈60-80\% & Q37 & 14-41\% & \textbf{27\%} \\
Q38∈60-80\% & Q39 & 21-46\% & \textbf{26\%} \\
Q42∈80-100\% & Q41 & 28-54\% & \textbf{25\%} \\
Q42∈20-40\% & Q45 & 57-84\% & \textbf{28\%} \\
Q44∈0-20\% & Q45 & 62-91\% & \textbf{29\%} \\
\bottomrule
\end{longtable}
\normalsize
\end{center}


窗口宽度 13-29\% 意味着只需扫描文章的 1/\allowbreak{}4 \textasciitilde{} 1/\allowbreak{}7 即可定位目标题。


\phantomsection
\subsection*{20.7 机械化决策树精要}
\addcontentsline{toc}{subsection}{20.7 机械化决策树精要}


\begin{CodeBlock}
Q36 ∈ 40-60% (n=19, 最常见):
  Q37 → 0-20% (58%)
  Q39 → 20-40% (63%)
  Q45 → 60-80% (42%)

Q36 ∈ 20-40% (n=21, 次常见):
  Q37 → 60-80% (52%)
  Q41 → 80-100% (48%)

双锚 Q36∈40-60% + Q38∈60-80% (n=8):
  Q37 → 0-20% (62%)
  Q39 → 20-40% (75%)

双锚 Q36∈20-40% + Q38∈40-60% (n=7):
  Q42 → 20-40% (71%)
  Q37 → 60-80% (57%)
  Q41 → 80-100% (57%)
\end{CodeBlock}


\phantomsection
\subsection*{20.8 与正文最终流程的对齐口径}
\addcontentsline{toc}{subsection}{20.8 与正文最终流程的对齐口径}


与第 14.1 节和第 17 章对齐后，统一流程应理解为：


\begin{enumerate}
\item Q36 起手（最稳定锚，σ=1.9）→\allowbreak{} 确定自己在哪条分支
\item 查决策树：Q36 带 →\allowbreak{} 推断 Q37/\allowbreak{}Q39 最可能区域
\item Q38 第二步（与 Q36 形成双锚，信息互补）
\item 用双锚决策树缩窗 →\allowbreak{} Q39 通常可缩至 20-40\%（75\% 命中）
\item Q40 或 Q42 第三步 →\allowbreak{} 三锚定位
\item 后续用弱排除和条件窗口处理剩余题
\item Q41 放最后（最难，但此时候选已很少）
\end{enumerate}


\textbf{不变铁律：}

\begin{itemize}
\item 相邻题永远不相邻（Q43/\allowbreak{}Q44 差>40\% 达 61\%）
\item 首段和末段是干扰高概率区（选率比期望低 12\%）
\item 做到第 7-8 题后，排除法价值上升；但已用字母不能绝对排除，因为同段重复合法
\end{itemize}


\phantomsection
\section*{21. 附录 G：工作包与图表入口索引}
\addcontentsline{toc}{section}{21. 附录 G：工作包与图表入口索引}


\phantomsection
\subsection*{21.1 图表文件入口}
\addcontentsline{toc}{subsection}{21.1 图表文件入口}


图表统一位于 \CodeInline{cet 6/\_analysis\_output/}：


\begin{itemize}
\item \CodeInline{question\_position\_boxplot.png}
\item \CodeInline{question\_position\_band\_distribution.png}
\item \CodeInline{question\_position\_band\_heatmap.png}
\item \CodeInline{question\_position\_curve.png}
\item \CodeInline{question\_answer\_distribution.png}
\item \CodeInline{letter\_to\_question\_heatmap.png}
\end{itemize}


\phantomsection
\subsection*{21.2 逐篇/\allowbreak{}逐批 Markdown 工作包入口}
\addcontentsline{toc}{subsection}{21.2 逐篇/逐批 Markdown 工作包入口}


以下目录仍然保留，承担“原始核对过程留痕”和“逐篇工作底稿”作用：


\begin{itemize}
\item \CodeInline{cet 6/\_analysis\_output/verification\_packets/}：逐篇人工核对包
\item \CodeInline{cet 6/\_analysis\_output/recheck\_packets/}：\CodeInline{sample.txt} 二次复核逐篇包
\item \CodeInline{cet 6/\_analysis\_output/recheck\_reviews/}：分批复核结论稿
\end{itemize}


如果阅读目标是“理解 Section B 规律、看最终结论、看考试打法”，优先看本主报告即可；如果阅读目标是“追溯某篇文章的逐题核对过程”，再进入上述目录查看对应 Markdown。

\end{document}
